\documentclass[output=paper,colorlinks,citecolor=brown]{langscibook}
\ChapterDOI{10.5281/zenodo.21237013}
\author{Georges-Jean Pinault\orcid{}\affiliation{École Pratique des Haute Études, Paris}}
%\ORCIDs{}

\title{Reassessing the emergence of the PIE *-\emph{eh\textsubscript{2}}-factitive present type}

\abstract{The PIE *-\textit{eh\textsubscript{2}}-factitive present type is currently known under the canonical instance of *\textit{néu̯-eh\textsubscript{2}}- ‘to make new’ (Hitt.\ \textit{nēwaḫḫ}-, Lat.\ \textit{(re)nouāre}, etc.), as based on the thematic adjective *\textit{néu̯-o}- ‘new’. The type has reflexes in Latin, Greek, Germanic, Balto-Slavic, Celtic, and possibly Tocharian, mostly in the thematic form \mbox{*-\textit{eh\textsubscript{2}-i̯e/o}-} (like  denominatives\is{denominative!verb} from abstract *-\textit{eh\textsubscript{2}}-stems). The Anatolian facts have confirmed that this formation was originally athematic and inflected according to the *\textit{h\textsubscript{2}e}-conjugation, reflected by the \textit{ḫi}-conjugation of the Hitt.\ \textit{nēwaḫḫ}-type, while the other Anatolian languages (Luwian, Lycian, Lydian) adopted the \textit{mi}-conjugation of the comparable types. According to the revision of the data by \citet{Sasseville2021}, the original function of this type in Anatolian was the derivation of denominal verbs from substantives, which was then extended to adjectives of various sorts. Several instances conclusively point to -\textit{ah\textsubscript{2}}-abstracts as the basis, which raises the issue of the so-called “\isi{conversion}” of a noun into an athematic present as the original PIE derivational process. As a preliminary investigation, the paper reviews other alleged reflexes of this process in other present types, which seem to point to denominative\is{denominative!verb} formations without apparent suffix, alternatively with zero suffix, for instance PIE *\textit{g\textsuperscript{u̯}ih\textsubscript{3}u̯é/ó}- ‘to live’, vis-à-vis the PIE adjective *\textit{g\textsuperscript{u̯}ih\textsubscript{3}u̯ó}- ‘living’. It will be shown that these apparent cases of \isi{conversion} of a noun into a verb are illusory and result from analogical developments, which presuppose the PIE existence of the “classical” full grade thematic present type, to wit *\textit{bʰér-e/o}-.} 


\IfFileExists{../localcommands.tex}{
   \addbibresource{../localbibliography.bib}
   \usepackage{tabularx,multicol}
\usepackage{url}
\urlstyle{same}

 
\usepackage{langsci-optional}
\usepackage{langsci-lgr}
\usepackage{langsci-gb4e}
\usepackage[linguistics,edges]{forest}
\usepackage{tikz-qtree}
\usetikzlibrary{positioning}
\usetikzlibrary{patterns.meta}
\usepackage{multirow}
\usepackage{pgfplots}
\usepackage{setspace}
\usepackage{amsmath,stackengine}
% \usepackage{fontspec}
% \usepackage{tipa} % TIPA is banned
\usepackage{langsci-textipa}
\usepackage{langsci-branding}
\usepackage{longtable}
\usepackage{tabto}

\usepackage{enumitem}

%\usepackage{lscape} %Querformat chapter 4

   % \newcommand*{\orcid}{}


\makeatletter
\let\thetitle\@title
\let\theauthor\@author
\makeatother

\newcommand{\togglepaper}[1][0]{
   \bibliography{../localbibliography}
   \papernote{\scriptsize\normalfont
     \theauthor.
     \titleTemp.
     To appear in:
    Laura Grestenberger, Viktoria Reiter \& Melanie Malzahn (eds.).
    \emph{Deadjectival verb formation in Indo-European and beyond}.
     Berlin: Language Science Press. [preliminary page numbering]
   }
   \pagenumbering{roman}
   \setcounter{chapter}{#1}
   \addtocounter{chapter}{-1}
}

\newbool{bookcompile}
\booltrue{bookcompile}
\newcommand{\bookorchapter}[2]{\ifbool{bookcompile}{#1}{#2}}

%Sonderzeichen
%idg palatales k
\newcommand{\kinvbreve}{k\raisebox{0.6ex}{\hspace{-0.05em}\char"0311}}

\newcommand{\jac}{%
  \textit{\char"0237}% italic dotless j (ȷ)
  \hspace{-0.01em}% fine-tune horizontal spacing
  \raisebox{0.01ex}{\char"0301}% raised combining acute
}

%Baltic circumflex l
\newcommand\ltilded{%
\hspace{-0.2em}%
  \stackon[-1.5pt]{l}{\smash{\kern0.2em\~{}}}%
  \hspace{-0.2em}%
}

%Slav. wedge+ double acute
\newcommand{\HighH}[1]{%
    \stackon[-1.2ex]{#1}{\kern0.1em\H{}\kern-0.1em}%
}

%Lith dot tilde VICKY diese b.-sl. Sonderzeichen BRINGEN MICH UM!!
\newcommand{\tildedot}[1]{%
    \ensurestackMath{%
        \stackon[-0.27ex]{% Tilde layer 
            \stackon[-0.15ex]{% Dot layer 
                #1%
            }{\hspace{0.1em}\cdot\hspace{-0.13em}}%
        }{\hspace{0.1em}\text{\~{}}\hspace{-0.13em}%
    }%
}}

%Lith dot acute
\newcommand{\dotacute}[1]{%
    \ensurestackMath{%
        \stackon[-1.1ex]{% acute layer 
            \stackon[-0.15ex]{% Dot layer 
                #1%
            }{\hspace{0.1em}\cdot\hspace{-0.12em}}%
        }{\hspace{0.1em}\text{\'{}}\hspace{-0.12em}%
    }%
}}



% \setlength{\emergencystretch}{2pt} 

 \pgfplotsset{compat=1.18}


\defcitealias{AbB_1}{AbB 1}
\defcitealias{AbB_7}{AbB 7}
\defcitealias{ABL_1}{ABL 1}
\defcitealias{Vries1977}{AEW}
\defcitealias{ALEW}{ALEW}
\defcitealias{AnGr}{AnGr}
\defcitealias{ARM_1}{ARM 1}
\defcitealias{BT}{BT}
\defcitealias{black_concise_2000}{CDA}
\defcitealias{CHD}{CHD}
\defcitealias{CT28}{CT 28}
\defcitealias{CT34}{CT 34}
\defcitealias{CT39}{CT 39}
\defcitealias{Cleasby1874}{CV}
\defcitealias{Beekes2010}{EDG}
\defcitealias{eDiAna}{eDiAna}
\defcitealias{eDIL}{eDIL}
\defcitealias{EDL}{EDL}
\defcitealias{Kroonen2013}{EDPG}
\defcitealias{ESJS}{ESJS}
\defcitealias{EVP}{EVP}
\defcitealias{EWA1988}{EWA}
\defcitealias{EWAI}{EWAia I}
\defcitealias{EWAII}{EWAia II}
\defcitealias{EWGP}{EWGP}
\defcitealias{GDLI1961}{GDLI}
\defcitealias{GPC}{GPC}
\defcitealias{GRADIT1999}{GRADIT}
\defcitealias{HW}{HW}
\defcitealias{IEW}{IEW}
\defcitealias{KAH_2}{KAH 2}
\defcitealias{KAR_1}{KAR 1}
\defcitealias{KAR_2}{KAR 2}
\defcitealias{KEWA}{KEWA}
\defcitealias{Labat_TDP}{Labat TDP}
\defcitealias{lambert_BWL}{Lambert BWL}
\defcitealias{LEIA}{LEIA}
\defcitealias{Rix2001}{LIV\textsuperscript{2}}
\defcitealias{LIVaddenda}{LIV\textsuperscript{2} {A}ddenda}
\defcitealias{LKG1971}{LKG}
\defcitealias{LKZ}{LKŽ}
\defcitealias{LP}{LP}
\defcitealias{MhdGr}{MhdGr}
\defcitealias{MLLVG1959}{MLLVG}
\defcitealias{NIL}{NIL}
\defcitealias{OgnS}{OgnS}
\defcitealias{ONP}{ONP}
\defcitealias{RIMA_2}{RIMA 2}
\defcitealias{TIM_2}{TIM 2}
\defcitealias{WAP}{WAP}
\defcitealias{YOS_2}{YOS 2}






\newindex{wan}{wdx}{wnd}{Form index}
\newcommand{\iw}[2]{\index[wan]{#1!{#2}\xspace}}
\newcommand{\iwi}[3][]{\index[wan]{#2!{#3}\xspace#1}{#3}}
\newcommand{\iwit}[3][]{\rule{0pt}{0pt}#1\index[wan]{#2!{#3}}{#3}}


\newindex{van}{vdx}{vnd}{Form index}
\newindex{ban}{bdx}{bnd}{Book index}

% \newcommand{\iw}[3][]{\index[wan]{#2!\textit{#3}}}


% \providecommand{\iwi}[3][]{\iw{#3}\textit{#3}}


\newcommand{\indexnote}[1]{{\footnotesize\sffamily\raggedright\normalfont\color{gray}#1}}

   %% hyphenation points for line breaks
%% Normally, automatic hyphenation in LaTeX is very good
%% If a word is mis-hyphenated, add it to this file
%%
%% add information to TeX file before \begin{document} with:
%% %% hyphenation points for line breaks
%% Normally, automatic hyphenation in LaTeX is very good
%% If a word is mis-hyphenated, add it to this file
%%
%% add information to TeX file before \begin{document} with:
%% %% hyphenation points for line breaks
%% Normally, automatic hyphenation in LaTeX is very good
%% If a word is mis-hyphenated, add it to this file
%%
%% add information to TeX file before \begin{document} with:
%% \include{localhyphenation}
\hyphenation{
    par-a-digm non-de-adj-ect-iv-al -pre-sent -pre-sents for-ma-tion-al in-do-ger-ma-ni-sche in-do-ger-ma-ni-schen Rei-chert Rei-ter dia-chrony ety-mo-lo-gi-sches Ost-ger-ma-ni-schen ger-ma-ni-schen alt-in-di-schen mor-pho-lo-gi-scher Alt-indo-ari-sch-en lexico-graphy Schrij-ver
}
% aktuell in Bibliographie: stud-ies, histor-ical Soci-età
\hyphenation{
    par-a-digm non-de-adj-ect-iv-al -pre-sent -pre-sents for-ma-tion-al in-do-ger-ma-ni-sche in-do-ger-ma-ni-schen Rei-chert Rei-ter dia-chrony ety-mo-lo-gi-sches Ost-ger-ma-ni-schen ger-ma-ni-schen alt-in-di-schen mor-pho-lo-gi-scher Alt-indo-ari-sch-en lexico-graphy Schrij-ver
}
% aktuell in Bibliographie: stud-ies, histor-ical Soci-età
\hyphenation{
    par-a-digm non-de-adj-ect-iv-al -pre-sent -pre-sents for-ma-tion-al in-do-ger-ma-ni-sche in-do-ger-ma-ni-schen Rei-chert Rei-ter dia-chrony ety-mo-lo-gi-sches Ost-ger-ma-ni-schen ger-ma-ni-schen alt-in-di-schen mor-pho-lo-gi-scher Alt-indo-ari-sch-en lexico-graphy Schrij-ver
}
% aktuell in Bibliographie: stud-ies, histor-ical Soci-età
   \boolfalse{bookcompile}
   \togglepaper[8]%%chapternumber
}{}

\begin{document}
\SetupAffiliations{mark style=none}
\maketitle


 
\section{Vedic \textit{rāj}- ‘to rule, be king’}
\subsection{Verb and corresponding nouns with a lengthened grade} 

The Vedic verb \iwi{Sanskrit}{\textit{rāj}-} (RV +) is currently glossed by an array of notions which leaves as somewhat uncertain the basic meaning, either ‘to rule, reign’, or ‘to shine, beam’, compare ‘herrschen, König sein, strahlen, glänzen’ (\citetalias[ 444]{EWAII}). This semantic ambiguity is actually linked to derivational and etymological issues. The corresponding root noun\is{root!noun} \textit{rā́j}-\iw{Sanskrit}{\textit{rā́j}- 'ruler'}  clearly means ‘king, ruler’. It is rare as independent noun (RV 3x), while being mostly found as second compound member,\is{compounding} in \textit{sam-rā́j}-, \iw{Sanskrit}{\textit{samrā́j}-} masc.\ (RV 38x, fem. \textit{sam-rā́jñī}- \iw{Sanskrit}{\textit{samrā́jñī}-} 4x), \textit{vi-rā́j}- \iw{Sanskrit}{\textit{virā́j}-} (7x), \textit{sva-rā́j}- \iw{Sanskrit}{\textit{svarā́j}-} (20x), \textit{jyeṣṭha-rā́j}- \iw{Sanskrit}{\textit{jyeṣṭharā́j}-} (2x), \textit{eka-rā́j}- \iw{Sanskrit}{\textit{ekarā́j}-} (1x), \textit{apna-rā́j}- \iw{Sanskrit}{\textit{apnarā́j}-} (1x), \textit{vane-rā́j}- \iw{Sanskrit}{\textit{vanerā́j}-} (1x), likewise in Old Iranian, see Av.\ \iw{Avestan}{\textit{bǝrǝzirāz}-} \textit{bǝrǝzi-rāz}-.\footnote{Survey of the RVic material by \citet[446--451]{Scarlata1999}.} Those compounds\is{compounding} belong to the type of verbal governing compound with root nouns\is{root!noun} as second member. The normal word for `king, ruler', is \iwi{Sanskrit}{\textit{rā́jan}-}, masc.\ (RV 256x), nom.sg.\ \textit{rā́jā}, acc.\ \textit{rā́jānam}, gen.\ \textit{rā́j}\textit{ñ}\textit{as}, dat.\ \textit{rā́j}\textit{ñ}\textit{e},  nom.pl.\ \textit{rā́jānas}, etc. If one adopts the current reconstruction of the root, this item can be traced back to a noun \iwi{PIE}{*\textit{h\textsubscript{3}rḗg̑-on}-}, with the suffix suggesting AK \is{amphikinetic}  inflection.\footnote{Henceforth, I use the following abbreviations for the standardly assumed PIE inflectional types: AS = \isi{acrostatic}, PK = \isi{proterokinetic}, HK = \isi{hysterokinetic}, AK = \isi{amphikinetic}, see \citet{Widmer2004}, \citet{Fortson2010}.}  At face value, it looks like an individualizing derivative (\iwi{PIE}{*-\textit{on}-/-\textit{n}-}) besides *\textit{h\textsubscript{3}rḗg̑}-,\iw{PIE}{*\textit{h\textsubscript{3}rḗg̑}-/*\textit{h\textsubscript{3}rég̑}- (noun)} which would be an abstract meaning ‘ruling, rulership’. This root noun\is{root!noun} is effectively attested as animate, and masculine, nom.sg.\ *\textit{h\textsubscript{3}rḗg̑-s}, which is reflected by RV \textit{rā́ṭ}, matching Lat.\  \iwi{Latin}{\textit{rēx}}, OIr.\ \iwi{OIr}{\textit{rí}}, Gaul.\ \iwi{Gaulish}{-\textit{rix}}, but as animate, it could also be feminine, without specific marking of the gender, as testified by RV 5.46.8b. According to \citet{Weiss2017}, PIE *\textit{h\textsubscript{3}rḗg̑}-\iw{PIE}{*\textit{h\textsubscript{3}rḗg̑}-/*\textit{h\textsubscript{3}rég̑}- (noun)} may have originally been a neuter root noun\is{root!noun} at the time when PIE had no feminine, meaning ‘ruling, directing’. The basis of the nasal stem was the locative sg.\ of this root noun:\is{root!noun}  *\textit{h\textsubscript{3}rēg̑-én} ‘in the power’ (> Ved.\  \textit{rāján-i}, \iw{Sanskrit}{\textit{rājáni}} RV 10.49.4c), from which the heteroclitic neuter stem found in Avestan, \textit{rāzarǝ}, gen.sg.\ \textit{rāzǝṇg}, instr.sg.\ \textit{rašnā}, gen.pl.\ \textit{rāšnąm} \iw{Avestan}{\textit{rāzarə̄̆}, \textit{rašn}-} was backformed.\is{backformation} The \isi{delocatival} animate yielded the nasal stem \iwi{PIE}{*\textit{h\textsubscript{3}rḗg̑-on}-} masc.\ ‘being in power’ > ‘ruler’, hence Ved.\  \iwi{Sanskrit}{\textit{rā́jan}-}. The -\textit{n}-stem abstract ‘rulership’ is also the basis of \iwi{PIE}{*\textit{h\textsubscript{3}rḗg̑-n-ih\textsubscript{2}}} fem.\ \isi{possessive} ‘having rulership’ (female) > ‘queen’, cf.\  Ved.\  \iwi{Sanskrit}{\textit{rā́jñī}-}, OIr.\ \iwi{OIr}{\textit{rigain}}, Lat.\ \iwi{Latin}{\textit{rēgīna}} < *\textit{rēgnīnā} (as per \citealt[293--295]{Pinault2014}).\footnote{Adopted with some minor differences by \citet[796--797, 799]{Weiss2017}.} 

This reconstruction leaves intact the fact of the apparent coexistence of the root noun\is{root!noun} Ved.\  \textit{rā́j}-\iw{Sanskrit}{\textit{rā́j}- 'ruler'}  and of an athematic root present\is{root!present} of the verb \iwi{Sanskrit}{\textit{rāj}-}, as reflected by Ved.\  (RV) 3sg.act.\ \iwi{Sanskrit}{\textit{rā́ṣṭi}}, which is unexpected from a synchronic point of view. Projected back into PIE, the verb *\textit{h\textsubscript{3}rḗg̑-ti} \iw{PIE}{*\textit{h\textsubscript{3}rḗg̑}- (verb)} has been interpreted as a denominative verb\is{denominative!verb} formed without suffix from the noun *\textit{h\textsubscript{3}rḗg̑}-\iw{PIE}{*\textit{h\textsubscript{3}rḗg̑}-/*\textit{h\textsubscript{3}rég̑}- (noun)} because of the meaning of the former: ‘to rule’ = ‘to be king’ (s.\ \citetalias[ III: 56]{KEWA} and \citetalias[ 444]{EWAII} with refs.). This is of course linked to the etymology of ‘king’ itself, which connects this noun to the PIE root \iwi{PIE}{*\textit{h\textsubscript{3}reg̑}-} (\citetalias[ 304--305]{Rix2001}) ‘gerade richten, ausstrecken’.\footnote{See Adams in \citealt[329--331]{MalloryAdams}, with previous literature. I leave aside the speculative and unsuccessful attempts to derive *\textit{h\textsubscript{3}rḗg̑}-\iw{PIE}{*\textit{h\textsubscript{3}rḗg̑}-/*\textit{h\textsubscript{3}rég̑}- (noun)} ‘king’ from a different root such as \iwi{PIE}{*\textit{reh\textsubscript{1}g̑}-} or \iwi{PIE}{*\textit{h\textsubscript{x}reh\textsubscript{1}g̑}-} (more specifically \iwi{PIE}{*\textit{h\textsubscript{2}}\textit{reh\textsubscript{1}g̑}-} through the connection with Gk.\ \iwi{AG}{ἀρήγειν}  ‘to help’, \iwi{AG}{ἀρήγων} ‘helper’) by several scholars in the past; see the literature in \citetalias[ III 50]{KEWA} and \citetalias[ 445--446]{EWAII}. In the Indian tradition itself, Skt.\  \iwi{Sanskrit}{\textit{rā́jan}-} has also been connected to roots other than \iwi{Sanskrit}{\textit{rāj}-}, see \citealt{Hara1969}.} 

In Indo-Iranian, only Ancient Vedic has the root athematic present \iwi{Sanskrit}{\textit{rā́ṣṭi}}.\is{root!present}  The root thematic present\is{thematic!present} is much better attested: RV \iwi{Sanskrit}{\textit{rā́jati}}, matched by Av.\  \mbox{\iwi{Avestan}{\textit{°rāzaiti}},} cf.\  \textit{vī-rāzaiti}\iw{Avestan}{\textit{vīrāzaiti}} ‘rules’ Yt.14.47 < Indo-Iranian \iwi{PIIr}{*\textit{rā́ȷ́a-ti}}. In Iranian the long vowel of the root has been extended to the \isi{verbal adjective}, Av.\ \iwi{Avestan}{\textit{rāšta}-} ‘directed, arranged’, OP \iwi{OP}{\textit{rāšta}-} ‘straight’, MP \iwi{MP}{\textit{rāšt}} ‘true’, etc.\footnote{See further \citealt[196--198]{Cheung2007}, under *\textit{Hraz}- ‘to draw a line, to direct’.} In the RV, the root thematic present\is{thematic!present} \mbox{\iwi{Sanskrit}{\textit{rā́jati}}} is the prevalent verbal stem (49x), subjunctive (2x), \isi{participle} active \iwi{Sanskrit}{\textit{rā́jant}-} (8x), beside inf.\ \iwi{Sanskrit}{\textit{rājáse}} (2x). There are only a few forms of other stems, aor.act.3pl.\ \iwi{Sanskrit}{\textit{arājiṣuḥ}} (1x),\footnote{See \citealt[221--222]{Narten1964}, under \textit{rāj}- ‘glänzen’.} \isi{intensive} pres.ptcp.act.\ \is{participle} \iwi{Sanskrit}{\textit{rā́rajat}} (1x). Projected back into PIE, this would point to a present \iw{PIE}{*\textit{h\textsubscript{3}rḗg̑-e/o}-} *\textit{h\textsubscript{3}rḗg̑-e-ti}, besides \iw{PIE}{*\textit{h\textsubscript{3}rég̑-e/o}-}  *\textit{h\textsubscript{3}rég̑-e-ti}, the normal root thematic present\is{thematic!present}, reflected by Gk.\ \iwi{AG}{ὀρέγω}, Lat.\ \iwi{Latin}{\textit{regō}}, -\textit{ere}, OIr.\  \iwi{OIr}{\textit{rigid}}, Go.\ \iwi{Goth}{\textit{rikan}}.  The whole issue has been treated by \citet{Strunk1987}, with further literature. Strunk’s own solution (\citeyear[389--390]{Strunk1987}) is as follows: The root \iwi{PIE}{*\textit{h\textsubscript{3}reg̑}-} had an \isi{acrostatic} (Narten)\is{Narten type} root present\is{root!present} with a 3sg.act.\ indicative *\textit{h\textsubscript{3}rḗg̑-ti}, subjunctive *\textit{h\textsubscript{3}rég̑-e-t(i)}. The thematic root present\is{thematic!present} found in most languages reflects this subjunctive. Indo-Iranian \mbox{\iwi{PIIr}{*\textit{rā́ȷ́a-ti}}} is explainable by generalization of the length of the root vowel to the subjunctive, and then to the injunctive and indicative. However, the RV subjunctive 1sg.act.\ \iwi{Sanskrit}{\textit{vi-rā́jāni}} (10.159.6d = 10.174.5d, in relatively late hymns) constitutes poor evidence for a subjunctive of the athematic root present\is{root!present} because it is most probably based on the thematic stem \textit{rā́ja}-\iw{Sanskrit}{\textit{rā́jati}} which is widespread already in the RV.  As far as the verb is concerned, this doctrine has been endorsed by \citetalias[ 305]{Rix2001}. As a minor and feeble alternative, the long grade of the verb would have its source in the nominal basis. Apparently, \citet{Strunk1987} took the long vowel grade of the Narten\is{Narten type} present and of the root noun\is{root!noun} ‘king’ as a coincidence, even though they are from the same root. Schindler’s oral communication, mentioned in a footnote (\citeyear[390, fn.\ 3]{Strunk1987}), would point, with question mark, to a root noun\is{root!noun} with \isi{acrostatic} root ablaut, *\textit{rḗg̑}-/*\textit{rég̑}- \iw{PIE}{*\textit{h\textsubscript{3}rḗg̑}-/*\textit{h\textsubscript{3}rég̑}- (noun)} (according to the notation used by Strunk), with further generalization of the lengthened grade, while the expected alternation is found in the heteroclitic neuter noun, OAv.\  \textit{rāzarə̄} \iw{Avestan}{\textit{rāzarə̄̆}, \textit{rašn}-}  ‘statute, order’, besides instr.sg.\ \textit{rašnā} (already quoted above) to wit *\textit{rḗg̑-r̥}/*\textit{rég̑-n}-,\iw{PIE}{*\textit{h\textsubscript{3}rḗg̑-r̥}/*\textit{h\textsubscript{3}rég̑-n}-} comparing (implicitly) the word for ‘liver’.  This idea was later adopted in the framework of “Narten systems”\is{Narten type} favored by Schindler in the late period of his research and teaching in Vienna. Schindler assumed that some PIE roots displayed “Narten”\is{Narten type} behavior (i.e., lengthened grade of the root in place of normal full grade *\textit{e} in the strong stem, besides full grade in place of expected zero grade), both in their verbal and in their nominal forms. This is reflected, for instance, in the treatment of the derivatives of the root \iwi{PIE}{*\textit{(h\textsubscript{3})reg̑}-} ‘lenken, herrschen’ by \citet[29--31]{Widmer1997}.  This included the reconstruction of the root noun\is{root!noun} as an original abstract noun and of the heteroclitic neuter as \iwi{PIE}{\textit{*(h\textsubscript{3})rḗg̑}-\textit{r̥}/\textit{*(h\textsubscript{3})rég̑-n}-}. In this discussion, the formation of the root present\is{root!present} Ved.\  \iwi{Sanskrit}{\textit{rā́ṣṭi}} was not addressed, however, possibly because it was considered as marginal. The issue is not substantially modified by the reconstruction of the root as \iwi{PIE}{*\textit{h\textsubscript{3}reg̑}-}, which imposed itself in the meantime. 

\subsection{Explanation of Ved.\  \textit{rā́ṣṭi} ‘reigns’ }

The coexistence of Ved.\  \iwi{Sanskrit}{\textit{rā́ṣṭi}}, \iwi{Sanskrit}{\textit{rā́jati}} ‘rules, reigns’ and of the noun (°)\textit{rā́j}-\iw{Sanskrit}{\textit{rā́j}- 'ruler'}  'ruler' ‘ruler’ is  puzzling because these verbs cannot be considered denominative presents\is{denominative!verb} in the normal sense. One would not deny that the prevalent use of RV \iwi{Sanskrit}{\textit{rā́ṣṭi}} and \iwi{Sanskrit}{\textit{rā́jati}} ‘to rule’, i.e., `to be king’ has been influenced by the meaning of (°)\textit{rā́j}-\iw{Sanskrit}{\textit{rā́j}- 'ruler'}  ‘ruler’, independent of whatever the original meaning of the root \iwi{Sanskrit}{\textit{rāj}-} was. An alternative strategy, which was favored in the early stages of Sanskrit lexicography, is to set up two roots on the basis of the translations of the occurrences in the most ancient text, the collection of hymns of the RV. There is however no need, \textit{pace} \citet[1156]{Grassmann}, to set up two homophonous verbal roots: 1.\ glänzen, strahlen (‘shine, be bright’) 2.\ herrschen, gebieten über (‘rule (over)’, governing the \isi{genitive}).  Both verbs would share the same present formation \mbox{\iwi{Sanskrit}{\textit{rā́jati}},} ptcp.\ \is{participle} \textit{rā́jant}-, so that this solution seems uneconomic. In most cases, the meaning ‘to shine, be bright’ is only inferred from the context, because the agent is a shining god (Agni, fire) or object (sun, liquid). The separation of these two meanings is actually superfluous, as can be shown by checking the majority of standard translations of the RV, which consistently situate the understanding of \iwi{Sanskrit}{\textit{rāj}-} under the notions of ruling and kingship of the agent of this verb.\footnote{See \citealt{GeldnerRV} and Renou, \citetalias{EVP}.}  In the most recent translation of the RV by \citet{JamisonBrereton2014}, one finds ‘rule (over)’, ‘reign’ in most occurrences. Alternative possibilities are considered in a minority of instances, and mostly between brackets: ‘rule [/shine]’ in 3.56.8b, 8.7.1c, 8.19.31d, 9.75.3d, 10.170.1d, ‘reign [shine forth]’ in 5.55.2b,  ‘shine [/rule]’ in 8.60.15d, ‘shine forth’ in 8.14.10c (for the drinks of \textit{soma,} subject of the aorist \iwi{Sanskrit}{\textit{arājiṣuḥ}}, hapax, with \isi{preverb} \iwi{Sanskrit}{\textit{ví}}), 9.5.1b, 9.5.3b, 9.61.18b, 10.189.3a, ‘radiates’ in 9.71.7d. This survey shows that one cannot absolutely exclude the interpretation of several occurrences of \iwi{Sanskrit}{\textit{rāj}-} as connected to the notion of shining (up) or beaming. This should not be taken to entail that ‘king’ itself (\textit{rā́j}-,\iw{Sanskrit}{\textit{rā́j}- 'ruler'}  \iwi{Sanskrit}{\textit{rā́jan}-}) originally meant ‘the bright one’, following a by now obsolete interpretation. The meaning ‘king, ruler’ is anchored in the inherited lexicon of Indo-Aryan and is supported without any doubt by the external comparanda. However, the semantic link between PIE *\textit{h\textsubscript{3}rḗg̑}-\iw{PIE}{*\textit{h\textsubscript{3}rḗg̑}-/*\textit{h\textsubscript{3}rég̑}- (noun)} and the root \iwi{PIE}{*\textit{h\textsubscript{3}reg̑}-}, for which there have been somewhat different approaches, poses yet another issue which would lead us beyond the limits of the present paper. For the purposes of my argument, it is sufficient to recall that this etymological connection is accepted by a wide majority of scholars.\footnote{As an aside, it should be recalled that the Anatolian data show that this etymon is limited to the languages of Core IE, so that *\textit{h\textsubscript{3}rḗg̑}-\iw{PIE}{*\textit{h\textsubscript{3}rḗg̑}-/*\textit{h\textsubscript{3}rég̑}- (noun)} cannot be deemed any longer as the unique and prototypical noun for ‘king’ in PIE.}   The semantic link from ‘stretch, direct’ to ‘rule’ is certainly inherited. This was probably based on the concept of the ruler whose power and control covers a vast domain, his realm, and extends in several directions. In Ancient Vedic, the transfer ‘ruling in every direction’ > ‘shining forth’ is due to the metaphorical association of the glory and supreme power of the king with the brightness of the sun and the radiance of the weapons of the warriors. In this area of expressions of symbolic notions, one may also note the favored use of the preverbs\is{preverb} \textit{ví} ‘apart, asunder, in different directions’ and \textit{sám} ‘together’ with the verb \iwi{Sanskrit}{\textit{rāj}-} and with its nominal derivatives. For the occurrences in the RV where the reference to ‘shining’ seems to be more convenient and more exact than ‘ruling over’, one may also invoke some rhyming effect with the verb \iwi{Sanskrit}{\textit{bhrāj}-} ‘to shine, beam, glitter’\footnote{\citetalias[ 279--280]{EWAII}; PIE root \iwi{PIE}{*\textit{bʰreh\textsubscript{1}}\textit{g̑}-} in \citetalias[ 92]{Rix2001}.} in some cases, which also forms a thematic root present,\is{thematic!present} RV middle \iwi{Sanskrit}{\textit{bhrā́jate}}, which is par excellence the verb used to describe the shining of the fire god Agni.  In a somewhat conservative way, \citet[267--271]{Goto1987} sets up two verbs (i.e., two thematic root presents)\is{thematic!present} \iwi{Sanskrit}{\textit{rā́jati}} (1.\ strahlen, 2.\ herrschen), basically following Grassmann, but admits that the two meanings are difficult to discriminate in most instances. Moreover, there is no link between these two meanings and different verbal formations nor etymologies. The whole picture is murky.\footnote{See also the discussion by \citet[445--446]{Scarlata1999}.} 

Returning now to the problem of derivation, the issue of \iwi{Sanskrit}{\textit{rā́ṣṭi}} and \iwi{Sanskrit}{\textit{rā́jati}} as anomalous denominatives\is{denominative!verb} is actually a pseudo-problem. The basis for the athematic root present\is{root!present} is small indeed, since it its only found in two occurrences in the  3sg.\ active, once in the indicative, and once in the injunctive: 

\begin{exe}
\ex RV 1.104.4b\\  \gll
\iwi{Sanskrit}{\textit{rā́ṣṭi}} \textit{śū́raḥ}\\
reign.\textsc{3sg.prs.ind.act} champion.\textsc{nom.sg.m}\\
\glt `The champion reigns.' 
\end{exe}

\begin{exe}
    \ex  RV 6.12.5cd\\ \gll 
\textit{sadyó} \textit{yáḥ} \textit{syandró} \textit{víṣito} {\textit{dhávīyān}  |}
\textit{r̥ṇó} \textit{ná} \textit{tāyúr} \textit{áti} \textit{dhán(u)vā} {\textit{rāṭ} ||}\\
{immediately} who.\textsc{nom.sg.m} {hasty}.\textsc{nom.sg.m} loosened.\textsc{nom.sg.m} hurrying.\textsc{comp.nom.sg.m}  {debtor}.\textsc{nom.sg.m} like thief.\textsc{nom.sg.m} across wasteland.\textsc{acc.pl.n} reign.\textsc{3sg.prs.inj.act} \iw{Sanskrit}{\textit{rā́ṣṭi}}\\
\glt `He [Agni] who immediately, when unloosened, streams ever faster, like a debtor (turned) thief he has headed straight across the waste places.' (transl.\ \citealt[786]{JamisonBrereton2014}). 
\end{exe}


By chance (or by a single poet’s choice?), the same hymn features one of the few (three in total) occurrences of the nom.sg.\ \textit{rā́ṭ}, 6.12.1a, at the same place in the pāda. Besides \textit{rā́ṭ}, the simplex noun \textit{rā́j}-\iw{Sanskrit}{\textit{rā́j}- 'ruler'}  has no other forms in the RV. Therefore, it seems better to explain this root athematic present \is{root!present} as an innovation within Ancient Vedic, which remained unsuccessful compared to \iwi{Sanskrit}{\textit{rā́jati}}, inherited from Indo-Iranian \iwi{PIIr}{*\textit{rā́ȷ́a-ti}}. One should find a way to explain both presents, without resorting to the notion of denominative\is{denominative!verb} formations. My scenario is as follows.  The Indo-Iranian thematic present\is{thematic!present} was inherited parallel to the formations attested in other IE languages and originally with normal full grade of the root: \iwi{PIIr}{*\textit{ráȷ́a-ti}} < *\textit{h\textsubscript{3}rég̑-e-ti},\iw{PIE}{*\textit{h\textsubscript{3}rég̑-e/o}-} ptcp.\ \is{participle} \iwi{PIIr}{*\textit{ráȷ́ant}-} < \iwi{PIE}{*\textit{h\textsubscript{3}rég̑-o-nt}-}. Those forms were remade as \iwi{PIIr}{*\textit{rā́ȷ́a-ti}}, ptcp.\ \is{participle} \iwi{PIIr}{*\textit{rā́ȷ́ant}-} under the influence of the nasal stem \iwi{PIIr}{*\textit{rā́ȷ́an}-} < \iwi{PIE}{*\textit{h\textsubscript{3}rḗg̑-on}-}, and of the second compound member\is{compounding} \iw{PIIr}{\textit{*rā́ȷ́}-} °\textit{rā́ȷ́}- < \iwi{PIE}{\textit{*°h\textsubscript{3}rḗg̑}-}. The two latter items had been independently inherited. The nasal stem\is{nasal present} was derived via a \isi{delocatival} derivation as outlined above and became firmly entrenched as the designation of the ruler at some stage. The process of reshaping the verbal stem hinges on two crucial points. The first point is the equivalence of the second member of a verbal governing compound \is{compounding} with a \isi{participle} and an agent noun, to wit ‘ruling over’. The second point is the crucial role of the present \isi{participle} as indicative of the structure of the corresponding present stem. This will also be relevant in the interpretation of the thematic present\is{thematic!present} \iw{PIE}{*\textit{g\textsuperscript{u̯}ih\textsubscript{3}-u̯é/ó}-} *\textit{g\textsuperscript{u̯}ih\textsubscript{3}u̯é-ti}, which we will discuss in the second section of this paper. 

The king’s function was correlated with the territory where his power was in force and where the clans (Ved.\  \iwi{Sanskrit}{\textit{víś}-}) under his authority reside, live and prosper in peace, a notion covered in Ancient Vedic by the root \iwi{Sanskrit}{\textit{kṣay}-/\textit{kṣi}-} ‘to dwell’ and its derivatives, see especially \iwi{Sanskrit}{\textit{kṣéma}-} masc. ‘peace, peaceful dwelling’. The Vedic term \iwi{Sanskrit}{\textit{rāṣṭrá}-} nt.\ (RV 10x) ‘realm, kingdom’ (basis of \iwi{Sanskrit}{\textit{rā́ṣṭrī}-} ‘ruler’ 3x) has been formed after the model of \iwi{Sanskrit}{\textit{kṣétra}-} nt.\ ‘field, ground, country’ (RV 20x, plus in several compounds),\is{compounding} the latter being inherited, cf.\  OAv.\  \iwi{Avestan}{\textit{šōiθra}-} ‘district, country’. The phrase RV \textit{kṣétrasya} \iwi{Sanskrit}{\textit{páti}-} is matched by Av.\  \textit{šōiθrahe} \iwi{Avestan}{\textit{paiti}-} ‘lord, master of the field’ = ‘king’. This lexical configuration yielded the causal force for the association between noun and verb. The verb \iwi{Sanskrit}{\textit{kṣay}-/\textit{kṣi}-} has a root athematic present \is{root!present} inherited from Indo-Iranian and from PIE, namely RV 3sg.act.\ \iwi{Sanskrit}{\textit{kṣéti}}, 3pl.\ \textit{kṣiyánti}, etc. (\citetalias[ 643]{Rix2001}, root \iwi{PIE}{*\textit{t\kinvbreve ei̯}-}). Based on these inherited lexemes, a proportion took shape: noun \iwi{Sanskrit}{\textit{kṣétra}-} : root present\is{root!present} \iwi{Sanskrit}{\textit{kṣéti}} :: \iwi{Sanskrit}{\textit{rāṣṭrá}-} : X = \iwi{Sanskrit}{\textit{rā́ṣṭi}}. 

A proportional analogy such as this suffices to account for the root athematic present \is{root!present} and de facto eliminates one of the favorite instances of an alleged denominative verb\is{denominative!verb} without suffix, i.e., of \isi{conversion} of a noun into a verb. One further advantage of the whole solution is to dispense with the otherwise interesting idea of a “Narten”\is{Narten type} root in the case of the root \iwi{PIE}{*\textit{h\textsubscript{3}reg̑}-}. I may recall from personal recollection that Schindler was forcefully opposed to the idea of denominative verbs\is{denominative!verb} without marking of derivation by a suffix. The preceding paragraphs should be understood as a first modest answer to this recurring problem. 

\section{The PIE thematic present *\emph{g\textsuperscript{u̯}ih\textsubscript{3}u̯é-ti}}
\subsection{A PIE *-\emph{u}-present}  


A PIE thematic present\is{thematic!present} can be safely reconstructed on the basis of several IE languages: Ved.\ \iw{Sanskrit}{\textit{jīvá}-} \textit{jī́vati}, YAv.\  \iw{Avestan}{\textit{juua}-} \textit{juuaiti}, Lat.\   \iw{Latin}{\textit{uīuere}} \textit{uīuit}, OPr.\ \iwi{OPr}{\textit{giwa}}, Latv.\ \iwi{Latv}{\textit{dzîvu}}, OCS \iwi{OCS}{\textit{živǫ}}. As for the root grade, Toch.\ B \iwi{TB}{\textit{śāw}-}, stem CToch.\ \iwi{CToch}{*\textit{śāw’ä}-} > Toch.\ B \iwi{TB}{*\textit{śāyä}-}, pres. act.\ 3sg.\ \iwi{TB}{\textit{śaiṃ}} (besides CToch.\ \iwi{CToch}{*\textit{śāwæ}-} in Toch.\ A act.\ 3pl.\ \iwi{TA}{\textit{śāweñc}}) is ambiguous, since \iwi{TA}{*\textit{śāw}-} may reflect \iwi{PIE}{*\textit{g\textsuperscript{u̯}i̯oh\textsubscript{3}u̯}-} as well as \iwi{PIE}{*\textit{g\textsuperscript{u̯}ih\textsubscript{3}u̯}-}.\footnote{See \citealt[916--919]{Malzahn2010} with previous literature.} Gk.\  \iwi{AG}{ζώω} has probably been reshaped with full grade on the basis of the root athematic aorist\is{root!aorist} 3sg.act.\ *\textit{g\textsuperscript{u̯}(i)i̯éh\textsubscript{3}-t} (cf.\  Gk.\  \iwi{AG}{ἐβίων}). In Vedic, the retraction of the accent to the root was due to the restructuring of the neo-root\is{root!secondary} \iwi{Sanskrit}{\textit{jīv}-} as a plain thematic present.\is{thematic!present} 

The mainstream account (e.g., \citetalias[ 215--216]{Rix2001}, \citealt[142]{Jasanoff2003}) assumes a PIE thematic present\is{thematic!present} which originated as a \isi{thematization} of an  athematic \isi{*-\textit{u}-present} \iwi{PIE}{*\textit{g\textsuperscript{u̯}i̯éh\textsubscript{3}-u}-/*\textit{g\textsuperscript{u̯}ih\textsubscript{3}-u}-´}.  This type followed the \isi{*\textit{h\textsubscript{2}e}-conjugation}  and hence had a  3sg.\ *\textit{g\textsuperscript{u̯}i̯éh\textsubscript{3}-u̯-e} that was the nucleus of the later thematic paradigm (\citealt{Jasanoff2003}, \citeyear{Jasanoff2023}). To this scenario one can object that there are no reflexes of an athematic \isi{*-\textit{u}-present} from the root \iwi{PIE}{*\textit{g\textsuperscript{u̯}i̯eh\textsubscript{3}}-}, contrary to the reflexes of other verbs: \iwi{PIE}{*\textit{terh\textsubscript{2}-u}-} ‘overcome, cross over’ (Hitt.\ \iwi{Hitt}{\textit{tarḫu}-}, 3sg.act.\ \textit{taruḫzi}, \textit{tarḫuzzi}, Ved.\  mid.\ \iwi{Sanskrit}{\textit{tarute}}, thematic\is{thematic!present} Ved.\ \iwi{Sanskrit}{\textit{tū́rvati}}, YAv.\ \iwi{Avestan}{\textit{tauruua}-}), \iwi{PIE}{*\textit{leh\textsubscript{2}-u}-} ‘pour’ (or \iwi{PIE}{*\textit{leh\textsubscript{3}-u}-}, as per \citealt{Melchert2011}), Hitt.\ \textit{laḫu}- (3sg.mid.\ OH \textit{laḫuwāri}, ptcp.\ \is{participle} \textit{laḫu(w)ānt}-),  \iw{Hitt}{\textit{laḫu}-, \textit{laḫuwāri}}   \iwi{PIE}{*\textit{u̯er-u}-} ‘ward off’, \iwi{PIE}{*\textit{u̯el-u}-} ‘turn’, etc.\footnote{See Jasanoff \citeyear[141--143]{Jasanoff2003}; \citeyear[63--66]{Jasanoff2023} for these examples and a few others. \textit{Pace} \citealt[66]{Jasanoff2023}, Toch.\ B \iwi{TB}{\textit{śaumo}} ‘man, human’ (A \iwi{TA}{\textit{śom}} ‘boy’) cannot be projected back to PIE \iwi{PIE}{*\textit{g\textsuperscript{u̯}i̯éh\textsubscript{3}u-mō(n)}} based on the suffixed\is{suffixation} root. This item ought to be explained as created at the CToch.\ stage, as per \citealt[119--120]{Pinault2021}.}   

\subsection{Questioning the theory of conversion}

An alternative theory that was entertained for a long time is that \iw{PIE}{*\textit{g\textsuperscript{u̯}ih\textsubscript{3}-u̯é/ó}-}  \mbox{*\textit{g\textsuperscript{u̯}ih\textsubscript{3}u̯é-ti}} was a denominative present\is{denominative!verb} based on the adjective \iw{PIE}{*\textit{g\textsuperscript{u̯}ih\textsubscript{3}-u̯ó}-}   *\textit{g\textsuperscript{u̯}}\textit{ih\textsubscript{3}u̯ó-s} (nom.sg.masc.) reconstructed for Core IE as well: Ved.\  \iwi{Sanskrit}{\textit{jīvá}-}, Av.\  \iwi{Avestan}{\textit{juua}-}, Lat.\   \iwi{Latin}{\textit{uīuus}}, OIr.\  \iwi{OIr}{\textit{biu}}, Go.\ \iwi{Goth}{\textit{qius}}, Lith.\ \iwi{Lith}{\textit{gývas}}, Latv.\ \iwi{Latv}{\textit{dzîvs}}, OCS \iwi{OCS}{\textit{živŭ}}.\footnote{See \citealt[55]{Tichy2004}, with part of the previous literature.} The infinitive Lat.\  \iwi{Latin}{\textit{uīuere}}, Ved.\  \iwi{Sanskrit}{\textit{jīváse}} reflects a stem \iwi{PIE}{*\textit{g\textsuperscript{u̯}ih\textsubscript{3}-u̯e/os}-} (underlying also Toch.\ B \iwi{TB}{\textit{śaiṣṣe}} ‘world’, A \iwi{TA}{\textit{śoṣi}} ‘people’)\footnote{As per Pinault \citeyear[234]{Pinault2008}, \citeyear[117]{Pinault2021}.} of the abstract based on this adjective. 

This pair of adjective and verb has been put into service by \citet[273--277]{VillanuevaSvensson2021}, claiming that the thematic present\is{thematic!present} issued from the thematic adjective by a process of “\isi{conversion}”, which may have been a regular derivational process at an early PIE stage:

\begin{exe}
   \ex adj.\ *\textit{g\textsuperscript{u̯}ih\textsubscript{3}-u̯é/ó}- → verb 3sg.\ *\textit{g\textsuperscript{u̯}ih\textsubscript{3}-u̯e-t(i)}, 3pl.\ \*\textit{g\textsuperscript{u̯}ih\textsubscript{3}-u̯o-nt(i)}. 
\end{exe}



According to Villanueva Svensson\ia{Villanueva Svensson, Miguel}, this pair is but a peculiar instance, which has been exceptionally preserved in Core IE, of the derivation by \isi{conversion} of a thematic present\is{thematic!present} from a thematic adjective (\citeyear[273--277]{VillanuevaSvensson2021}), which was a purely nominal formation, not based on an alleged, but actually non-existent \isi{*-\textit{u}-present} (with strong stem *\textit{g\textsuperscript{u̯}i̯éh\textsubscript{3}-u}-).\iw{PIE}{*\textit{g\textsuperscript{u̯}i̯éh\textsubscript{3}-u}-/*\textit{g\textsuperscript{u̯}ih\textsubscript{3}-u}-´} This pair reflected a standard process which is hypothetically assumed for early Core IE, through which the simple thematic present\is{thematic!present} of the type \iwi{PIE}{*\textit{bʰér-e/o}-} ‘to carry’ was derived from a reconstructed \isi{verbal adjective} \iwi{PIE}{*\textit{bʰer-ó}-} ‘carrying’ (\citealt[277--284]{VillanuevaSvensson2021}). Consequently, the whole root thematic present\is{thematic!present} type emerged \textit{after} the separation of Anatolian from PIE and became widely productive\is{productivity} in Core IE. The latter point is the major claim of his paper, because Villanueva Svensson\ia{Villanueva Svensson, Miguel} dismisses previous theories about the place of the deverbative type\is{deverbal/deverbative!verb} \iwi{PIE}{*\textit{bʰér-e/o}-} in the PIE verbal system (\citeyear[264--273]{VillanuevaSvensson2021}).

The same process is then found in other pairs consisting of a thematic adjective and a thematic root present,\is{thematic!present} e.g. Gk.\  \iwi{AG}{λευκός} ‘white’ (\iwi{PIE}{*\textit{leu̯k-ó}-}): Ved.\  \iwi{Sanskrit}{\textit{rócate}} ‘shines’ \iw{PIE}{*\textit{léu̯k-e/o}-} (*\textit{léu̯k-e-to+i}), yielding the thematic root present \is{thematic!present} through additional adjustments of accentuation to obtain an intermediate \iwi{PIE}{*\textit{léu̯k-o}-} ‘shining’ that could serve as the basis for  \isi{conversion}. However, Villanueva Svensson\ia{Villanueva Svensson, Miguel} fails to account for the full grade and the root accentuation of the \iwi{PIE}{*\textit{bʰér-e/o}-} present type as well as the typical *-\textit{e}/\textit{o}-alternation of the thematic vowel in the inflection, except through circular and \textit{ad hoc} hypotheses (such as assuming that the original accentuation of the adjective \iwi{PIE}{*\textit{leu̯k-ó}-} > Gk.\  \iwi{AG}{λευκός} was **\textit{léu̯k-o}\nobreakdash-, just for complying with the accent of the present stem \iwi{PIE}{*\textit{léu̯k-e/o}-}, cf.\  \citealt[281--282, and fn.\ 43]{VillanuevaSvensson2021}).

An earlier variant of the \isi{conversion} account for the verb ‘living’ has been formulated by \citet[55]{Tichy2004} plainly as follows: “3.Sg.\ Ind.\ Präs.\ uridg.\ \iw{PIE}{*\textit{g\textsuperscript{u̯}ih\textsubscript{3}-u̯é/ó}-} *\textit{g\textsuperscript{u̯}ih\textsubscript{3}-u̯e-ti} ‘ist *\textit{g\textsuperscript{u̯}ih\textsubscript{3}u̯ó}-, lebt’.” This is however a mere pedagogical paraphrase instead of an explanation, which leaves out the real nominal form that could be used as a predicate in a nominal sentence, the nom.sg.\ \iw{PIE}{*\textit{g\textsuperscript{u̯}ih\textsubscript{3}-u̯ó}-}  *\textit{g\textsuperscript{u̯}ih\textsubscript{3}u̯ó-s}, in favor of the transposition of the thematic stem. Tichy quotes \citet[79]{Rix1994}, who pushes back the phenomenon into PIE: “hocharchaische Denominativbildung mit Nullsuffix”.\is{denominative!verb}

The basic problem is: What does “\isi{conversion}” really mean? In other words, how would it work in syntactic and morphological terms? Does it imply that PIE was not a fully-fledged inflectional language?  Under this account, there is no clearcut link through reinterpretation of a predicate thematic stem, with nominal endings, nom.sg.\ \textit{*°o-s} (animate) or \textit{*°o-m} (neuter) as verbal 3sg.act.\ \textit{*-e} (active or \isi{proto-middle}; active \textit{*°e-t}, middle \textit{*°e-to} (injunctive). How could this gap be bridged? This cannot be understood remotely in the spirit of “\isi{Watkins’ Law}” (\citealt{Watkins1969}), according to which a form *\textit{bʰér-e} with “\isi{proto-middle}” 3sg.\ ending was reinterpreted as a thematic stem with zero ending, triggering the generalization of the stem  \iw{PIE}{*\textit{bʰér-e/o}-} *\textit{bʰér-e}-.\footnote{See \citealt[224--227]{Jasanoff2003}. Watkins’ “law”\is{Watkins’ Law} was used as powerful tool for interpreting various verbal formations, see \citealt{Joseph1980} and \citealt{Janse2009} with further references.} 

The *\textit{°e} form of the thematic vowel did not surface in the nominal inflection except for the vocative and possibly in some adverbial forms. Nobody has claimed that the thematic inflection started from the imperative, while taking for granted that the imperative with zero ending in the second singular active could be akin to a vocative (and vice versa). Nor is there any basis for setting up a predicative adverb **\textit{bʰer-é} ‘while carrying’, which is mentioned here only for the sake of the argument. Therefore, it appears that the theory of \isi{conversion} cannot be pushed very far, since it entails many additional speculations.  

\citet[74--77]{Jasanoff2023}, on the other hand, argues that PIE did possess the full grade thematic present\is{thematic!present} type, which did not arise through \isi{conversion} but through a process of \isi{reanalysis} of the “proto-thematic” 3sg.\ active ending *-\textit{e} of the \isi{*\textit{h\textsubscript{2}e}-conjugation}, coexisting with active (\isi{transitive}) 1sg.\ *-\textit{e-h\textsubscript{2}e}, 2sg.\  *-\textit{e-th\textsubscript{2}e}, later recharacterized as injunctive 3sg.act.\ *-\textit{e-t}, hence indicative *-\textit{e-ti}. There are instances of standard root thematic presents\is{thematic!present} already in Anatolian, not to speak of the numerous complex suffixed thematic presents\is{thematic!present}, in \iwi{PIE}{\textit{*-i̯e/o}-}, \iwi{PIE}{*-\textit{s\kinvbreve e/o}-}, *-\textit{ei̯e/o}-\iw{PIE}{*-\textit{éi̯e/o}-} with active thematic inflection. As for the above pair itself, the scenario developed by \citet[62--71]{Jasanoff2023} would take the two items as independent formations from the \isi{*-\textit{u}-present}  which followed the \isi{*\textit{h\textsubscript{2}e}-conjugation} (type \iw{PIE}{*\textit{térh\textsubscript{2}-u̯}-} *\textit{térh\textsubscript{2}-u̯\nobreakdash-e}, with -\textit{i} added to the ending in the active, as per \citealt[141--143]{Jasanoff2003}): The present \iw{PIE}{*\textit{g\textsuperscript{u̯}ih\textsubscript{3}-u̯é/ó}-}  *\textit{g\textsuperscript{u̯}ih\textsubscript{3}-u̯e/o}- arose through \isi{thematization} based on the 3sg.act.\ *\textit{g\textsuperscript{u̯}ih\textsubscript{3}-u̯-e(+i)} (for expected **\textit{g\textsuperscript{u̯}i̯éh\textsubscript{3}-u̯-e}, presupposing levelling of the root ablaut) and was characteristically active, while the adjective was presumably a kind of thematic \isi{verbal adjective} based on the regular weak stem *\textit{g\textsuperscript{u̯}ih\textsubscript{3}-u̯}-. \iw{PIE}{*\textit{g\textsuperscript{u̯}i̯éh\textsubscript{3}-u}-/*\textit{g\textsuperscript{u̯}ih\textsubscript{3}-u}-´} The latter point remains somewhat in the dark, if not taken for granted. 

\subsection{An alternative derivation with a verbal adjective as intermediate item}
\label{sec:altder}

In the following, I will try a different course, based on well-known facts of nominal derivation and semantics. PIE \iwi{PIE}{*\textit{g\textsuperscript{u̯}ih\textsubscript{3}-u̯ó}-} ‘living’ was a derivative from the root \iwi{PIE}{*\textit{g\textsuperscript{u̯}i̯eh\textsubscript{3}}-} by means of the archaic verbal adjective-forming suffix\is{verbal adjective} \iwi{PIE}{*-\textit{u̯o}-},\footnote{This is more cogent than resorting to the oppositional suffix found in \iwi{PIE}{*\textit{de\kinvbreve s(i)-u̯o}-} ‘right’ (Gk.\  \iwi{AG}{δεξι(ϝ)ός}, Go.\ \iwi{Goth}{\textit{taihswo}}), since the basis of this term was not a verbal root.} which can be compared to its semantic opposite \iwi{PIE}{*\textit{mr̥-u̯ó}-} ‘dead’ > OIr.\  \iwi{OIr}{\textit{marb}} from the root \iwi{PIE}{*\textit{mer}-} ‘to die’, replaced by \iwi{PIE}{*\textit{mr̥-tó}-} > Ved.\  \iwi{Sanskrit}{\textit{mr̥tá}-}, Av.\  \iwi{Avestan}{\textit{mərəta}-}, Arm.\  \iwi{Armenian}{\textit{mard}}, etc., with (parallel) contamination into \iwi{PIE}{*\textit{mr̥-tu̯o}-} > Lat.\   \iwi{Latin}{\textit{mortuus}}, OCS \iwi{OCS}{\textit{mĭrtvŭ}}. The development of \iwi{PIE}{*\textit{g\textsuperscript{u̯}ih\textsubscript{3}-u̯ó}-} should therefore be viewed in the context of the semantic field which covers life, old age and death. The root \iwi{PIE}{*\textit{g̑erh\textsubscript{2}}-} ‘to wear out, age, make look older’ (\citetalias[ 165]{Rix2001}) formed a plain thematic \isi{verbal adjective} PIE \iwi{PIE}{*\textit{g̑erh\textsubscript{2}-ó}-} ‘decayed’ (Arm.\  \iwi{Armenian}{\textit{cer}}) and a parallel athematic synonymous \mbox{\iwi{PIE}{*\textit{g̑érh\textsubscript{2}-ont}-/*\textit{g̑r̥h\textsubscript{2}-n̥t}}-´} > Ved.\  \iwi{Sanskrit}{\textit{járant}-} and \iwi{Sanskrit}{\textit{jurat}-} (dat.sg.\ \textit{juraté}, gen.pl.\ \textit{juratā́m}), fem.\ \iwi{Sanskrit}{\textit{járatī}-}, Gk.\  \iwi{AG}{γέρων}, γέρoντ-, Oss.\ \iwi{Oss}{\textit{zærond}}) ‘old, decayed’, featuring AK \is{amphikinetic} inflection, and belonging to a Caland system\is{Caland!system (CS)}  (cf.\  the neuter abstract \iwi{PIE}{*\textit{g̑érh\textsubscript{2}-s}} > Gk.\  \iwi{AG}{γέρας} shifted to ‘honor, perk’ and remade as \iwi{AG}{γῆρας} under the influence of the aorist \iwi{AG}{ἔγηρᾱ}, as per \citealt{Stueber2002}: 84).\footnote{Consequently, the routine reconstruction of an \isi{acrostatic} neuter *-\textit{s}-stem, as per Widmer (\citeyear[20]{Widmer1997}, \citeyear[50]{Widmer2004}) and other scholars, is not necessary.}  The idea of *\textit{g̑érh\textsubscript{2}-ont}-\iw{PIE}{*\textit{g̑érh\textsubscript{2}-ont}-/*\textit{g̑r̥h\textsubscript{2}-n̥t}-´} as the \isi{participle} of a Narten\is{Narten type} verb, present or aorist, as per \citealt[21]{Widmer1997} (see also \citetalias[ 165, n.\ 2]{Rix2001}) does not withstand scrutiny. The Ancient Vedic verbal system included the \isi{transitive}-\isi{factitive} present \iwi{Sanskrit}{\textit{járati}} ‘to make old’, besides \isi{intransitive} stems such as the present RV \iwi{Sanskrit}{\textit{jū́ryati}}/AV \iwi{Sanskrit}{\textit{jī́ryati}} ‘to grow old, become frail’ and the \isi{stative} \isi{perfect} RV \isi{participle} act.\ \textit{jujurvā́n}, \textit{jujurúṣ}-, AV \textit{jajā́ra} ‘ist gealtert’. One cannot reconstruct any PIE aorist of this verb. An \textit{*-(o)n}-stem is reflected by the Toch.\ B stem \iwi{TB}{\textit{śrānä}-} masc. ‘adult man’,\footnote{\citealt[441, 484]{Pinault2008}, differently \citealt[705]{Adams2013}.}  nom.pl.\ \textit{śrāy}, obl.pl.\ \textit{śrānäṃ}, gen.pl.\ \textit{śrānäṃts}; \textit{śrāy} < *\textit{śärāyä} < \iwi{CToch}{*\textit{śärāñä}} < *\textit{g̑érh\textsubscript{2}-ōn-es} (alternatively, but less likely, *\textit{g̑érh\textsubscript{2}-n-es}), presupposing as stem \iwi{PIE}{*\textit{g̑érh\textsubscript{2}-ōn}-} < *\textit{g̑érh\textsubscript{2}-o-on}- ‘the old one’, with an individualizing suffix based on the thematic adjective mentioned above. 

After the model of *\textit{g̑érh\textsubscript{2}-ont}-,\iw{PIE}{*\textit{g̑érh\textsubscript{2}-ont}-/*\textit{g̑r̥h\textsubscript{2}-n̥t}-´} reanalyzed as *\textit{g̑érh\textsubscript{2}-o-nt}-, parallel to the synonymous \iwi{PIE}{*\textit{g̑erh\textsubscript{2}-ó}-}, \iwi{PIE}{*\textit{g̑érh\textsubscript{2}-ōn}-}, a remodelled form \iw{PIE}{*\textit{g\textsuperscript{u̯}ih\textsubscript{3}-u̯ó-nt}-} *\textit{g\textsuperscript{u̯}ih\textsubscript{3}u̯ó-nt}- meaning ‘the living one’ was created, opposed to \iwi{PIE}{*\textit{mér-to}-} ‘the dead one’ (Ved.\  \iwi{Sanskrit}{\textit{márta}-} masc.) ‘mortal, human’ and \iwi{PIE}{*\textit{mertōn}-} < *\textit{merto-on}- (Av.\  \iwi{Avestan}{\textit{marətān}-}, masc. ‘human’, nom.sg.\ \textit{marətā}, nom.pl.\ \textit{marətānō}). At a later stage, when the \isi{verbal adjective} formation in \iwi{PIE}{\textit{*-u̯ó}-} had been replaced by *-\textit{tó}- \iw{PIE}{*-\textit{to}-} and \iwi{PIE}{\textit{*-(e/o)nt}-},  \textit{g\textsuperscript{u̯}ih\textsubscript{3}u̯ónt}- could be reanalyzed as \iw{PIE}{*\textit{g\textsuperscript{u̯}ih\textsubscript{3}-u̯ó-nt}-} *\textit{g\textsuperscript{u̯}ih\textsubscript{3}u̯ó-nt}-, the \isi{participle} of a thematic verbal stem, whence the present injunctive act.\ 3pl.\ \iw{PIE}{*\textit{g\textsuperscript{u̯}ih\textsubscript{3}-u̯é/ó}-}  *\textit{g\textsuperscript{u̯}ih\textsubscript{3}u̯o-nt}, 3sg.\ *\textit{g\textsuperscript{u̯}ih\textsubscript{3}u̯e-t}.  The final stage of this process presupposes the PIE existence of root thematic presents\is{thematic!present} of the “classical” type, to wit \iwi{PIE}{*\textit{bʰér-e/o}-}, injunctive act.\ 3sg.\ *\textit{bʰér-e-t}, 3pl.\ *\textit{bʰér-o-nt}, \isi{participle} \iwi{PIE}{*\textit{bʰér-o-nt}-}.  If not for this precise root, one may surmise that this type existed in PIE, see, e.g., the \textit{*-i̯e/o}- and *-\textit{s\kinvbreve e/o}-presents with “complex” thematic suffixes.\is{*-\textit{i̯e/o}-present} \is{*-\textit{s\kinvbreve e/o}-present} The PIE existence of full grade thematic presents\is{thematic!present} of the type 3sg.act.\ *\textit{bʰér-e-ti} has recently been supported by new independent arguments,\footnote{See \citealt[72--78]{Jasanoff2023}.} even though the presents of this type were not as abundant at the early stage as they were later in Core IE.  

To conclude, the thematic present\is{thematic!present} \iw{PIE}{*\textit{g\textsuperscript{u̯}ih\textsubscript{3}-u̯é/ó}-} *\textit{g\textsuperscript{u̯}ih\textsubscript{3}u̯é/ó}- was wholly analogical, while being based ultimately on an adjectival stem. This derivation offers an additional instance of creating the inflectional forms of a finite present stem based on a \isi{verbal adjective} reinterpreted as \isi{participle}. My scenario thus differs markedly from the recent accounts of the pair \iwi{PIE}{*\textit{g\textsuperscript{u̯}ih\textsubscript{3}-u̯ó}-} ‘alive’ : *\textit{g\textsuperscript{u̯}ih\textsubscript{3}-u̯e/o}- \iw{PIE}{*\textit{g\textsuperscript{u̯}ih\textsubscript{3}-u̯é/ó}-} ‘to live’ by \citet{VillanuevaSvensson2021} and \citet{Jasanoff2023}. 

\section{PIE *\textit{néu̯-ah\textsubscript{2}}- ‘to make new’}
\subsection{Brief review of the Anatolian evidence} \label{sec:3.1}

The type Hitt.\ \iwi{Hitt}{\textit{nēwaḫḫ}-}\is{\textit{nēwaḫḫi}-type} ‘to make new’ has been much discussed in the literature.\footnote{Cf.\  \citealt[155, 240, 247, 253, 455--458]{Oettinger1979}, \citealt[132]{Melchert1997}, \citealt[139--141]{Jasanoff2003}. These references are not intended to be exhaustive. See also the introduction to this volume for an overview.} The Anatolian evidence has recently been reassessed by Sasseville (\citeyear{Sasseville2015} and \citeyear[16--77]{Sasseville2021}). The type\is{\textit{nēwaḫḫi}-type} was productive\is{productivity} in Hittite in forming deadjectival \is{deadjectival!verb} \is{factitive}  factitives  from thematic  and other kinds of stems: \iwi{Hitt}{\textit{nēwaḫḫ}-} ‘to make new’ : \iwi{Hitt}{\textit{nēwa}-} ‘new’, \iwi{Hitt}{\textit{maršaḫḫ}-} : \iwi{Hitt}{\textit{marša}-} ‘unholy’, \iwi{Hitt}{\textit{dannattaḫḫ}-} : \iwi{Hitt}{\textit{dannatta}-} ‘empty’, \iwi{Hitt}{\textit{ḫappinaḫḫ}-} : \iwi{Hitt}{\textit{ḫappina}-}, \iwi{Hitt}{\textit{ḫappinant}-} ‘rich’, \iwi{Hitt}{\textit{šuppiyaḫḫ}-} : \iwi{Hitt}{\textit{šuppi}-} ‘pure’, \iwi{Hitt}{\textit{tepawaḫḫ}-} ‘to make \mbox{little’ :} \iwi{Hitt}{\textit{tēpu}-} ‘little’; from substantives, e.g. \iwi{Hitt}{\textit{armaḫḫ}-} ‘to make pregnant’ < ‘make round like the moon’ : \iwi{Hitt}{\textit{arma}-} c.\ ‘moon’, \iwi{Hitt}{\textit{luriyaḫḫ}-} ‘to humiliate’ : \iwi{Hitt}{\textit{luri}-} c.\ ‘humiliation’, \iwi{Hitt}{\textit{šagiyaḫḫ}-} ‘to give a sign’ : \iwi{Hitt}{\textit{šagāi}-} c.\ ‘sign’. This type belongs to the derivational set of deadjectival verbs\is{deadjectival!verb} together with \isi{fientive} -\textit{ēšš}-\iw{Hitt}{-\textit{ešš}-} and \isi{factitive} \iwi{Hitt}{-\textit{nu}-} verbs. 

The \textit{°aḫḫ}- type followed the \textit{ḫi-}conjugation in Hittite (OH and MH), which was later replaced by the \isi{\textit{mi}-conjugation}.  By contrast, in the other Anatolian languages, the \isi{\textit{mi}-conjugation} prevails, with non-lenited endings in the corresponding presents in Luwian /-\textit{a}-/, \iw{CLuw}{-\textit{a}-}  \iw{HLuw}{-\textit{a}-} Lycian \iwi{Lycian}{-\textit{a}-} and Lydian \iwi{Lydian}{-\textit{a}-}. Those verbs are mostly derived from substantives, rarely from adjectives: HLuw.\  \iwi{HLuw}{\textit{tabariya}-} ‘to command’ : \iwi{HLuw}{\textit{tabariya}-} c.\ ‘command’, Lyc.\  A \iwi{Lycian}{\textit{prñnawa}-} ‘to build (a grave-\mbox{house)’ :} \iwi{Lycian}{\textit{prñnawa}-} c.\ ‘grave-house’, Lyc.\  B \iwi{Lycian}{\textit{muwa}-} ‘to strengthen’ : coll. \textit{muwa} ‘might, power’, Lyc.\  B \iwi{Lycian}{\textit{qla}-} ‘to build a precinct’ : \iwi{Lycian}{\textit{qala}-/\textit{qla}- c.} ‘precinct’, Lyc.\  A \iwi{Lycian}{\textit{sm͂ma}-} ‘to bind, oblige’ : \iwi{Lycian}{\textit{sm͂ma}-} c.\ ‘binding, obligation’, Lyd.\ \iwi{Lydian}{\textit{pita}-} ‘to make a donation’ (‘to make pay’ according to Garnier, p.c.) : *\textit{pita}- c.\ ‘donation’, cf.\  Lyc.\  A \iwi{Lycian}{\textit{pijata}-} c.\ ‘gift’, Hitt.\ \iwi{Hitt}{\textit{piyetta-, pitta}-}, coll. ‘allotment, payment’. They could also be derived from adjectives or agent nouns: CLuw.\ \iwi{CLuw}{\textit{tannatta}-} ‘to make empty, devastate’ : \iwi{CLuw}{\textit{tannatta/i}-} ‘empty’, CLuw.\ \iwi{CLuw}{\textit{arkammanalla}-} ‘to make someone into a tributary’ : \iwi{CLuw}{\textit{arkammanalla/i}-} ‘tribute-bearing’, Lyc.\ A \iwi{Lycian}{\textit{kumaza}-} ‘to perform a sacrifice’ < ‘to act as a priest’ : \iwi{Lycian}{\textit{kumaza/i}-} ‘sacrificing priest’, HLuw.\  \iwi{HLuw}{\textit{hantawatta}-} ‘to act as a \mbox{king’ :} \iwi{HLuw}{\textit{hantawatt(i)-} c.} ‘king’. Those agent nouns actually reflect formations in \textit{°ā} < \iwi{PIE}{*-\textit{eh\textsubscript{2}}-}, the individualizing suffix, e.g., Lyc.\  \iwi{Lycian}{*-\textit{za}-} < \iwi{PIE}{*-\textit{ti̯ah\textsubscript{2}}-}, derived from secondary adjectives\is{derivation!stem-based} in \iwi{PIE}{*-\textit{ti̯o}-} (see \citealt{Sasseville2014,Sasseville2018}).\footnote{Following \citealt{Melchert2014}, see further below.} 

\subsection{Considering the issue of conversion}

The usual comparanda for the type are the *-\textit{ā}-\isi{factitive}s in Latin, Greek, Germanic, Balto-Slavic, and Celtic. This is often summarized by the equation Hitt.\ \iwi{Hitt}{\textit{nēwaḫḫ}-},\is{\textit{nēwaḫḫi}-type} Lat.\ \iwi{Latin}{\textit{(re)nouāre}}, Gk.\ \iwi{AG}{νεάω} ‘to plough up anew’, OHG \iwi{OHG}{\textit{niuwōn}}, despite some secondary restructurings. Thus, Gk.\  \iwi{AG}{νεάω} is actually associated with \iwi{AG}{νειός} fem.\ ‘new-ploughed land’ (< *\textit{neu̯ii̯ó}-), not with the adjective \iwi{AG}{νέος} ‘new’ and OHG \iwi{OHG}{\textit{niuwōn}} is based on the alternative Germanic adjective ‘new’, OHG \iwi{OHG}{\textit{niuwi}} (< \iwi{PIE}{*\textit{neu̯(i)i̯o}-}, cf.\  Ved.\  \iwi{Sanskrit}{\textit{náv(i)ya}-}, \isi{doublet} of \iwi{Sanskrit}{\textit{náva}-}). Nonetheless, the derivational process can be summarized by the following scheme in PIE:

\begin{exe}
    \ex thematic adjective \iwi{PIE}{*\textit{néu̯-o}-} ‘new’ →\\ verb *\textit{néu̯-e-h\textsubscript{2}}-\iw{PIE}{*\textit{néu̯-e(-)h\textsubscript{2}}- ‘to make new’} ‘to make new’, 3sg.act.\ *\textit{néu̯-e-h\textsubscript{2}-e(i)}, 3pl.act.\ *\textit{néu̯-e-h\textsubscript{2}-n̥ti} 
\end{exe}



This entails that the thematic vowel of the base took on the *-\textit{e}-grade, assuming that the verbal suffix is reconstructed as plain *-\textit{h\textsubscript{2}}-. Alternatively, the suffix \iwi{PIE}{*-\textit{eh\textsubscript{2}}-} would be substituted for the final \iwi{PIE}{*-\textit{o}-} of the base. This moot point is ordinarily not considered in the scholarship, even though it has significant consequences for ascertaining the origin of the formation. 

This present formation was subsequently thematized\is{thematization} to \iwi{PIE}{*-\textit{eh\textsubscript{2}-i̯e/o}-}, which is found in Latin, Greek, Balto-Slavic, Vedic, and possibly Tocharian,\footnote{\citealt[579]{Pinault2008}, further discussion in \citealt[393--402]{Malzahn2010}.}  as denominatives \is{denominative!verb} from reflexes of PIE *-\textit{ah\textsubscript{2}}-stems. This type became quite productive\is{productivity} in several languages. One may leave aside henceforth the case of the alternative \isi{thematization} in \iwi{PIE}{*-\textit{eh\textsubscript{2}-e/o}-}, which became productive\is{productivity} in Balto-Slavic.\footnote{See Jasanoff \citeyear[140--141]{Jasanoff2003}; \citeyear[200--201]{Jasanoff2017BSl};  \citealt[393--395]{VillanuevaSvensson2020}.} This type has left reflexes in Aeolic and Mycenaean Greek contract verbs, pointing to an originally athematic 3sg.act.\ *\textit{neu̯ah\textsubscript{2}ei̯}, cf.\  Lesb.\ {-αι}
\iw{AG}{-αι (3sg.)}
as 3sg.act.\ in the athematic paradigm of contract verbs (with 1sg.act.\ \iwi{AG}{-ᾱμι} instead of \iwi{AG}{-άω}), and Myc.\ \textit{te-re-ja}, 3sg.act.\ present < *\textit{telehi̯ai̯}, from the verb \iwi{PG}{*\textit{telehi̯ā̆}- ‘to accomplish, fulfill one’s obligations’} \citep{Rau2009myc}. The basis can be either the adjective *\textit{telehi̯o}-, a syncopated form of *\textit{telehii̯o}- \iw{PG}{*\textit{teleh(i)i̯o}-} > Gk.\ \iwi{AG}{τέλειος} ‘perfect, accomplished’ (derived itself from the neuter \textit{s}-stem reflected by \iwi{AG}{τέλoς} ‘service, obligation’), or an abstract \mbox{\iwi{PG}{*\textit{telehi̯ā}}} < \iwi{PIE}{*\textit{teles-i̯ah\textsubscript{2}}} (\citealt[181, fn.\ 4]{Rau2009myc}). The Myc.\ infinitive \textit{te-re-ja-e} \iw{MycG}{\textit{te-re-ja}, \textit{te-re-ja-e}} < \iwi{PG}{\textit{*°āhen}} (cf.\  Lesb.\ \iwi{AG}{-ᾱν}) has the expected form for an athematic stem, since the ending may reflect ultimately a virtual locative sg.\ \iwi{PIE}{\textit{*°ah\textsubscript{2}-sen}}, similar to other enlargements of abstracts, such as the well-known types \iwi{PIE}{*\textit{-ah\textsubscript{2}-ter/n}-}, \iwi{PIE}{*-\textit{ah\textsubscript{2}-u̯er/n}-}.\footnote{\citealt[32--34, 135, 198, fn.\ 4]{Nussbaum1986}; \citealt[347--354, 380--382, 414--417]{Rieken1999}.}

The verbal suffix *-\textit{h\textsubscript{2}}- (after thematic stems) with \isi{factitive} value has not been accounted for from the etymological point of view. There is no comparable suffix in the verbal derivational system\is{derivation!verbal} of PIE. The segmentation itself (as \iwi{PIE}{*-\textit{eh\textsubscript{2}}-} or *-\textit{e-h\textsubscript{2}}- if the basis was a thematic adjective) remains elusive. If one assumes the “\isi{conversion}” of a nominal stem in \iwi{PIE}{*\textit{-(e)h\textsubscript{2}}-} as the basis of a verbal stem, it would be required to ascertain what the underlying nominal derivative was, e.g., an abstract/collective neuter in \iwi{PIE}{*-\textit{eh\textsubscript{2}}-} (*-\textit{e-h\textsubscript{2}}-) or an individualizing animate in \iwi{PIE}{*-\textit{eh\textsubscript{2}}-}.\footnote{Cf.\  \citealt{Melchert2014}, \citealt[274--275, 306]{Nussbaum2014coll}, and \citealt{FellnerGrestenberger2016}   on \mbox{*\textit{-(e)h\textsubscript{2}}-} as marker of substantivization.}  This issue would not matter so much for scholars (such as one anonymous reviewer of the present paper) who consider “that at an early stage of Indo-European nominal and verbal inflection differed only in the endings, but not in the stem formation”, so that the suffix \iwi{PIE}{*-\textit{eh\textsubscript{2}}-} could form both nominal and verbal stems. This hypothesis is based on a very simplistic and glottogonic view of the morphological reconstruction. In any case, it would not account for the \isi{*\textit{h\textsubscript{2}e}-conjugation} of the PIE *\textit{néu̯-ah\textsubscript{2}}- present type, nor for its \isi{factitive} meaning. The present stems in \iwi{PIE}{*-\textit{i}-} and \iwi{PIE}{*-\textit{u}-} \is{*-\textit{u}-present} cannot be straightforwardly identified with suffixed nominal stems either, which feature several inflectional patterns which are not all found in the verbs. In addition, this approach ignores the existence of denominative presents\is{denominative!verb} in PIE and in the IE daughter languages, as well as of verbal adjectives\is{verbal adjective} provided with specific and different morphemes. I will therefore not consider it further.

As shown by Sasseville (\citeyear[287--292]{Sasseville2014}; \citeyear[69--70]{Sasseville2021}), the original Anatolian function of verbal -\textit{aḫḫ-} was the derivation of verbs from substantives, which was then extended to adjectives, as in the Hitt.\ \textit{nēwaḫḫ}-type.\is{\textit{nēwaḫḫi}-type} Several instances conclusively point to *-\textit{ah\textsubscript{2}}-abstracts as the base, for instance Lyd. \iwi{Lydian}{\textit{pita}-} ‘to make a donation’ from a *-\textit{tah\textsubscript{2}}-stem, Hitt.\ \iwi{Hitt}{\textit{waštaḫḫ}-} ‘to sin, to offend’ vs.\ \iwi{Hitt}{\textit{waštai}-} c.\ ‘sin’, CLuw. \iwi{CLuw}{\textit{wašta}-} nt.\ ‘sin’, Hitt.\ \iwi{Hitt}{\textit{armaḫḫ}-} ‘to make round like the moon, to impregnate’ vs.\ \iwi{Hitt}{\textit{arma}-} c., CLuw., HLuw.\ \textit{arma}- c.\  ‘moon’, \iw{CLuw}{\textit{arma}-} \iw{HLuw}{\textit{arma}-} Lyc.\  A \iwi{Lycian}{\textit{arm͂ma}-} c.\ ‘moon-god’ < stem in \iwi{PIE}{*-\textit{mah\textsubscript{2}}-}. Concerning the latter example, I follow the derivation proposed by Sasseville (\citeyear[282--284]{Sasseville2014}; \citeyear[69--71]{Sasseville2021}), which  complies with the expected semantic evolution of a verb of the \isi{factitive} Hitt.\ \textit{°aḫḫ}-type. Therefore, this account is certainly to be preferred to the one suggested by one anonymous reviewer for the original meaning of this verb: ‘to bring to the months’, or perhaps ‘to bring to the month (of birth)’. This remains arbitrary, and at variance with the \isi{factitive} meaning warranted by other instances of \textit{°aḫḫ}-verbs which are clearly based on substantives. There is no support for the alleged directive (or allative) meaning which would be taken by the basis of this denominative formation.\is{denominative!verb} 

In addition to these \isi{factitive} verbs which were apparently derived without any overt suffix, Anatolian had normal denominatives\is{denominative!verb} formed with the suffix \iwi{PIE}{*-\textit{i̯e/o}-},\is{*-\textit{i̯e/o}-present} featuring \isi{essive} value: e.g. Hitt.\ \iwi{Hitt}{\textit{armā(i)}-} < *°\textit{ah\textsubscript{2}-i̯e/o}- \iw{PIE}{*-\textit{ah\textsubscript{2}-i̯e/o}-} ‘to be round like the moon’ > ‘be pregnant’, contrasting with \iwi{Hitt}{\textit{armaḫḫ}-} ‘to make round like the moon’, quoted above.\footnote{As argued cogently by \citet[287--288]{Sasseville2015}.}  

\citet[292--293]{Sasseville2015} assumes the “\isi{conversion}” of *-\textit{ah\textsubscript{2}}-abstracts into \isi{factitive} *-\textit{ah\textsubscript{2}}-verbs, challenging the standard view of the \textit{nēwaḫḫ}-type\is{\textit{nēwaḫḫi}-type} as derived from thematic adjectives, for Anatolian as well as for PIE. The derivation from thematic adjectives would be a development specific to Hittite within Anatolian and independently of Core IE. The process of restructuring of a nominal suffix into a verbal suffix is taken for granted, without setting up any intermediate step. This approach is endorsed by Villanueva Svensson\ia{Villanueva Svensson, Miguel} (\citeyear[391--392]{VillanuevaSvensson2020}, \citeyear[275--276]{VillanuevaSvensson2021}) because it provides welcome support to his own account which argues for the derivation of the full grade thematic present\is{thematic!present} type from thematic verbal adjectives.\is{verbal adjective}  This entails that there was originally no proper verbal suffix in verbs with a 3sg.act.\ *-\textit{ah\textsubscript{2}-e}, which followed the \isi{*\textit{h\textsubscript{2}e}-conjugation}. One may surmise that the thematic renewal as \iwi{PIE}{*-\textit{ah\textsubscript{2}-e/o}-} found in Balto-Slavic issued originally from the reinterpretation of this 3sg.\ active form. The missing link could be a further individualizing derivative in *-\textit{on}-\iw{PIE}{*-\textit{on}-/-\textit{n}-} or \iwi{PIE}{*-\textit{(e/o)nt}-}, the latter being reinterpreted as a \isi{participle}. This would not explain the original *-\textit{h\textsubscript{2}}\textit{e} conjugation, however. Moreover, such secondary derivatives\is{derivation!stem-based} of abstracts or animate substantives in \iwi{PIE}{*-\textit{ah\textsubscript{2}}-} did not leave any reflex in unmistakable form. The case of adj. \iwi{PIE}{*\textit{g\textsuperscript{u̯}ih\textsubscript{3}-u̯ó}-} ‘living’ vs.\ verb \iw{PIE}{*\textit{g\textsuperscript{u̯}ih\textsubscript{3}-u̯é/ó}-} \*\textit{g\textsuperscript{u̯}ih\textsubscript{3}-u̯é-ti} ‘(s)he lives, is alive’ does not add support to the “\isi{conversion}” account because this present \textit{does} feature the normal thematic conjugation\is{thematic!present} of the *\textit{bʰér-e-ti} type; see Section \ref{sec:altder} above. 

\subsection{The Vedic evidence}

In previous discussions of possible comparanda of the Hitt.\ \textit{nēwaḫḫ}-type,\is{\textit{nēwaḫḫi}-type} Indo-Iranian has not been mentioned so far. However, some intriguing evidence in Ancient Vedic can be detected: The present RV  \iwi{Sanskrit}{\textit{vr̥ṣāyáti/-te} ‘to make rain’} (3x), which must be distinguished from \iw{Sanskrit}{\textit{vr̥ṣāyáte} `to behave like a bull'} \textit{vr̥ṣāyáte} (RV 7x, ptcp.\ \is{participle} \textit{vr̥ṣāyámāṇa}- 2x) ‘to behave like a bull, act the bull’, denominative\is{denominative!verb}  from \iwi{Sanskrit}{\textit{vŕ̥ṣan}-} masc.\ (see also the parallel adjective \iwi{Sanskrit}{\textit{vr̥ṣayú}-} ‘acting the bull’ in the herd among the cows, RV 9.77.5d). The structure of the former item has not yet been cogently accounted for. The reading of the occurrences does not leave any doubt: 

\begin{exe}
    \ex RV 10.98.1d [to Br̥haspati]\\
    \gll \textit{sá} \textit{parjányaṃ} \textit{śáṃtanave} {\textit{vr̥ṣāya} \iw{Sanskrit}{\textit{vr̥ṣāyáti/-te} ‘to make rain’} ||}\\
    \textsc{dem.pron.nom.sg.m} P.\textsc{acc.sg} S.\textsc{dat.sg} rain.make.\textsc{2sg.imp.act}\\
    \glt `Make Parjanya rain for Śaṃtanu.' (in a plea for rain by Devāpi, priest and poet. Transl.\ \citealt[1555]{JamisonBrereton2014})
\end{exe}

\begin{exe}
    \ex RV 9.71.3ab\\ \gll \textit{ádribhiḥ} \textit{sutáḥ} \textit{pavate} \textit{gábhast(i)yor}, {\textit{vr̥ṣāyáte} \iw{Sanskrit}{\textit{vr̥ṣāyáti/-te} ‘to make rain’}} \textit{nábhasā} \textit{vépate} {\textit{matī́} |} \\
    stone.\textsc{ins.pl} pressed.\textsc{nom.sg.m} purify.\textsc{3sg.prs.ind.mid} hand.\textsc{loc.du} rain.make.\textsc{3sg.prs.ind.mid} cloud.\textsc{ins.sg} tremble.\textsc{3sg.prs.ind.mid} thought.\textsc{ins.sg} \\
    \glt `Pressed by the stones, he (Soma) purifies himself between the two hands. With his cloud he makes rain (the rain = his juice); he trembles (in poetic inspiration) with his thought.' (differently from \citealt[1304]{JamisonBrereton2014}, who translate ‘he acts the bull [/rains]’).\footnote{The first option corresponds to Geldner’s choice: “durch die Regenwolke wird er wie ein Bulle, durch die Dichtung wird er beredt” (\citealt[III, 63]{GeldnerRV}). See further discussion by Renou (\citetalias{EVP} VIII 83 and IX 81), who notes that the reference to the rain is implied by \textit{nábhas}- as complement, evoking the cloud.} 
\end{exe}


\begin{exe}
    \ex RV 10.44.4ab\\ \gll \textit{evā́} \textit{pátiṃ} \textit{droṇa-sā́caṃ} \textit{sácetasam}, \textit{ūrjá} \textit{skambháṃ} \textit{dharúṇa} \textit{ā́} {\textit{vr̥ṣāyase} \iw{Sanskrit}{\textit{vr̥ṣāyáti/-te} ‘to make rain’} |} \\
    thus lord.\textsc{acc.sg} cup-companion.\textsc{acc.sg} like.minded.\textsc{acc.sg} nourishment.\textsc{gen.sg} prop.\textsc{acc.sg} support.\textsc{loc.sg} \textsc{prvb} rain.make.\textsc{2sg.prs.ind.mid}\\
    \glt `Just in this way you [Indra] drench yourself in the lord [= soma], the companion of the cup, the like-minded one, the prop of nourishment upon its support.' (transl.\ \citealt[1449]{JamisonBrereton2014}).
\end{exe}



This present stem provided the basis for the secondary -\textit{iṣ}-aorist\is{derivation!stem-based}  \iwi{Sanskrit}{\textit{ā́vr̥ṣāyiṣata}} (VS II.31), following the pres.\ imperative  \textit{ā́vr̥ṣāyadhvam} (cf.\  \citealt[62, 292]{Narten1964}), both with \isi{preverb} \textit{ā́}, as in RV 10.44.4b. The formation of the present is not discussed by Narten\ia{Narten, Johanna}.  

This present \textit{vr̥ṣāyáti/-te} \iw{Sanskrit}{\textit{vr̥ṣāyáti/-te} ‘to make rain’} can hardly represent a formation of the \textit{gr̥bhāyáti} type, based originally on nasal infixed present,\is{nasal present} to wit \iwi{Sanskrit}{\textit{gr̥bhṇā́ti}} from \iwi{Sanskrit}{\textit{grabh\textsuperscript{i}}-} ‘to seize’, with \iwi{Sanskrit}{-\textit{āyá}-} < *-\textit{aH-yá}- < *-\textit{n̥H-i̯é/o}-,\footnote{\citealt[107]{Goto2013}.  The type is reconstructed as \iwi{PIE}{*-\textit{n̥h\textsubscript{2}-i̯e/o}-} by \citet[123--124]{Jasanoff2003}. The root \iwi{Sanskrit}{\textit{grabh\textsuperscript{i}}-} is currently traced back to \iwi{PIE}{*\textit{gʰrebh\textsubscript{2}}-} (\citetalias[ 201]{Rix2001}), but the precise laryngeal in this present type does not matter in the present context.}  since \iwi{Sanskrit}{\textit{varṣ}-/\textit{vr̥ṣ}-} (\textit{aniṭ} root) did not form a -\textit{nā-}present.\is{nasal present}  Nor can it be a plain \isi{essive} or \isi{fientive} denominative\is{denominative!verb}  from a noun \iwi{Sanskrit}{*\textit{vr̥ṣā}-} (or \iwi{Sanskrit}{*\textit{vr̥ṣa}-}) meaning ‘rain’, because it has \isi{factitive} meaning.\footnote{\textit{Pace} \citealt[131]{Goto2013}. In any case, one should renounce searching the origin of the \isi{factitive} in the denominative\is{denominative!verb}  of ‘bull’, through the metaphorical assimilation of the jet of sperm by the male to the flow of rain. This cannot be supported by the problematic etymological interpretation of PIE ‘bull’ itself, from a root that would mean ‘to impregnate’, \textit{vel sim}.}  This verb did not survive long, since it was in \isi{competition} with the regular \isi{causative} \iwi{Sanskrit}{\textit{varṣáyati}}  ‘to make rain’ (RV 3x, \isi{participle} \textit{varṣáyant}- 1x, with complement \textit{dyā́m} ‘sky’ or \textit{vr̥ṣṭím} ‘rain’), besides the \isi{intransitive} pres. \iwi{Sanskrit}{\textit{várṣati}} ‘to rain’.\footnote{\citealt[117--118]{Jamison1983}.} 

Instead, I propose that \textit{vr̥ṣā-yá-ti} \iw{Sanskrit}{\textit{vr̥ṣāyáti/-te} ‘to make rain’} replaced a \isi{factitive} athematic present \textit{*u̯r̥šā}- < *\textit{h\textsubscript{2}u̯r̥seh\textsubscript{2}}-. The latter cannot be based on an adjective \iwi{PIE}{*\textit{h\textsubscript{2}u̯r̥s-ó}-} ‘rainy, raining’, which is unattested. The basic nominal derivatives of the PIE root \iwi{PIE}{*\textit{h\textsubscript{2}u̯ers}-} are well known (see \citealt[255--256]{Nussbaum2017}): \isi{agentive}s which have been variously substantivized, namely \iwi{PIE}{*\textit{h\textsubscript{2}u̯órs-o}-} or \iwi{PIE}{*\textit{h\textsubscript{2}u̯ors-ó}-} (Hitt.\ \iwi{Hitt}{\textit{warša}-} c.), \iwi{PIE}{*\textit{h\textsubscript{2}u̯ors-ó}-} or \iwi{PIE}{*\textit{h\textsubscript{2}u̯ers-ó}-} (Ved.\  \iwi{Sanskrit}{\textit{varṣá}-} nt.), collective *\textit{h\textsubscript{2}u̯ersé-h\textsubscript{2}} (Ved.\  \textit{varṣā́} ‘the rains’, ‘rainy season’), abstract \iw{PIE}{*\textit{h\textsubscript{2}u̯érseh\textsubscript{2}}-} *\textit{h\textsubscript{2}u̯érse-h\textsubscript{2}}- (Gk.\  \iwi{AG}{ἐέρση, ἀέρση}  ‘dew’). The potential stem \iwi{PIE}{*\textit{h\textsubscript{2}u̯r̥s-éh\textsubscript{2}}-} could be a primary derivative\is{derivation!deradical} of the type Gk.\  \iwi{AG}{φυγή}, Lat.\   \iwi{Latin}{\textit{fuga}} ‘flight’ or a replacement of the expected root noun\is{root!noun} \iwi{PIE}{*\textit{h\textsubscript{2}u̯érs}-/*\textit{h\textsubscript{2}u̯r̥s}-} ‘raining’ reflected by RV \iwi{Sanskrit}{\textit{prāvr̥ṣ}-} fem.\ (2x) ‘beginning of the rainy season’, derived adj.\ \iwi{Sanskrit}{\textit{prāvr̥ṣī́ṇa}-} (1x).\footnote{\citealt[525]{Scarlata1999}.} 

In order to explain the \isi{factitive} present from an abstract, the best way lies in the syntactic reconstruction of an intermediary nominal sentence: 

\begin{exe}
  \ex  \isi{dative} sg.\ in predicative function\is{predicative!function} *\textit{h\textsubscript{2}u̯r̥séh\textsubscript{2}-ei} ‘(s)he/it (is) for the raining’ > ‘(s)he/it makes rain’. 
\end{exe}


Next, the sequence could be resegmented as a 3sg.\ present *\textit{h\textsubscript{2}u̯r̥séh\textsubscript{2}-e+i} with a primary ending, hence as a verbal stem *\textit{h\textsubscript{2}u̯r̥séh\textsubscript{2}}-,\iw{PIE}{*\textit{h\textsubscript{2}u̯r̥s-éh\textsubscript{2}}-} which was later thematized\is{thematization} as Indo-Iranian \iwi{PIIr}{*\textit{u̯r̥šā-i̯á-ti}}, featuring the accent on the thematic suffix as in regular deverbative\is{deverbal/deverbative!verb} or denominative presents.\is{denominative!verb} As for the syntax, one can compare the well-known construction of the \isi{dative} sg.\ of verbal nouns in predicative function,\is{predicative!function} such as Vedic RV \textit{ná dábhāya} ‘not to be deceived’ (5.44.2c, 7.91.2a, 9.73.8a) or \textit{ná (…) ādábhe} (RV 8.21.16d) with the same meaning. These negative phrases in nominal sentences are equivalent in the RV to the privative adjectives \iwi{Sanskrit}{\textit{ádabdha}-} (48x), \iwi{Sanskrit}{\textit{ádābh(i)ya}-} (34x) ‘undeceivable’, \iwi{Sanskrit}{\textit{adábha}-} (5.86.5b) ‘undeceptive’.\footnote{On the relationship of infinitives and quasi-infinitives with verbal adjectives,\is{verbal adjective} see \citealt{Renou1937}.}  Compare also the reading of the final \isi{dative} of a root noun\is{root!noun} used as quasi-infinitive with \isi{factitive} meaning in RV 1.89.5c \textit{védasām ásat vr̥dhé} ‘he will be for increasing/growing of possessions’ = ‘for letting grow the possessions’.\footnote{See \citealt[95, 96]{Gippert1978} and \citealt[23--25]{HettrichStueber2018} for further instances.} 

\subsection{Consequences for Anatolian and PIE}

The derivational scheme devised for explaining Ved.\  \textit{vr̥ṣāyáti} \iw{Sanskrit}{\textit{vr̥ṣāyáti/-te} ‘to make rain’} applies in principle to any *-\textit{eh\textsubscript{2}}- abstract \isi{dative} sg.\ *\textit{X-eh\textsubscript{2}-ei̯}. \iw{PIE}{*\textit{-eh\textsubscript{2}-ei̯}} This was the precise basis of the so-called \isi{conversion}. The predicative \textit{*[X-eh\textsubscript{2}]-ei̯} in nominal sentences was reinterpreted as \iw{PIE}{*\textit{-eh\textsubscript{2}-ei̯}} \textit{X-eh\textsubscript{2}-ei̯}, thus as 3sg.\ active present in *-\textit{e+i} of the \isi{*\textit{h\textsubscript{2}e}-conjugation}, with *\textit{X-eh\textsubscript{2}-} reanalyzed as a denominative verbal stem.\is{denominative!verb} This happened after the 3sg.\ ending of the \isi{*\textit{h\textsubscript{2}e}-conjugation} was reanalyzed as an active ending, which yielded the Anatolian \isi{\textit{ḫi}-conjugation}. The further step to be made consists in transposing this process into PIE to generate the *-\textit{eh\textsubscript{2}}-\isi{factitive} present type. This can be supported by a concrete example found in Anatolian: \iwi{PAn}{*\textit{armaH-ei̯}}  ‘(in)to the moon’ {$\approx$} ‘to become/make like the moon’ > Hitt.\ 3sg.act.\ \textit{armaḫḫ-i} \iw{Hitt}{\textit{armaḫḫ}-} ‘makes N like the moon’ = ‘makes N pregnant’, see Section \ref{sec:3.1} above. This is consistent with the prevalent derivation of this \isi{factitive} present type from abstract substantives, and by extension from other *-\textit{ah\textsubscript{2}}- substantives, as shown by Sasseville (\citeyear{Sasseville2021}, see again Section  \ref{sec:3.1}). The form \iwi{PAn}{*\textit{armaH-ei̯}} ‘(in)to the moon’ could in principle generate a labile verb,\is{labile verbs} ‘to become/make (N) [round] like the moon’. But it was used transitively,\is{transitive} in contrast to the \isi{essive}/\isi{fientive} suffixed\is{suffixation} denominative\is{denominative!verb} Hitt.\ \iwi{Hitt}{\textit{armā(i)}-} ‘to be pregnant’ < \iwi{PAn}{*\textit{°ah\textsubscript{2}-i̯e/o}-} ‘to be(come) [round] like the moon’. Hence, the former predicative was verbalized as Hitt.\ 3sg.act.\ \textit{armaḫḫ-i} \iw{Hitt}{\textit{armaḫḫ}-} ‘makes N [round] like the moon’ > ‘makes N pregnant’.  One should assume that the \isi{transitive} reading prevailed for *-\textit{ah\textsubscript{2}}-abstracts used predicatively\is{predicative!usage} in the \isi{dative} singular with an overt nominal complement.  

On the other hand, one cannot deny that the \isi{productivity} of the type \iwi{PIE}{*\textit{neu̯eh\textsubscript{2}}-} ‘to make new’ in Core IE as based on thematic adjectives had a primeval source in PIE, although its prevalence and extension to other types of adjectives is a phenomenon proper to Hittite in the Anatolian language family. One may surmise that this development was only embryonic in PIE. Nonetheless, it is likely that some *-\textit{eh\textsubscript{2}}- abstracts could be derived from thematic adjectives. Therefore, the chain of derivation was extended with respect to the base, as in ex.\ (\ref{ex:9}).  

\newpage
\begin{exe}
   \ex adjective \iwi{PIE}{*\textit{néu̯-o}-}  
→\\ abstract *\textit{néu̯-e-h\textsubscript{2}}-\iw{PIE}{*\textit{néu̯-eh\textsubscript{2}}- `newness, news'} ‘newness’ →\\ \isi{dative} sg.\ *\textit{néu̯-eh\textsubscript{2}-ei̯} ‘for the purpose of newness, into newness’ 
→\\ verb *\textit{néu̯-eh\textsubscript{2}-e+i}\iw{PIE}{*\textit{néu̯-e(-)h\textsubscript{2}}- ‘to make new’} ‘(s)he makes X into newness’ > ‘(s)he makes X new’. \label{ex:9}
\end{exe}

This derivational chain was then reduced to \iwi{PIE}{*\textit{néu̯-o}-} → *\textit{néu̯eh\textsubscript{2}-e+i}, resegmented as *\textit{néu̯-e-h\textsubscript{2}-e+i},\iw{PIE}{*\textit{néu̯-e(-)h\textsubscript{2}}- ‘to make new’} in the deadjectival \is{deadjectival!verb} Hitt.\ \textit{nēwaḫḫ}-type\is{\textit{nēwaḫḫi}-type} and its cognates in other IE languages. The restructuring of a complex chain by “jumping over” the intermediate link is a well-known phenomenon. The surface effect of the process was finally the apparent substitution of the nominal thematic vowel by the verbal suffix, thus contrasting with the *-\textit{o-i̯é/ó}- or *-\textit{e-i̯é/ó}- denominative presents.\is{denominative!verb} 

This “de-datival” \is{dative} derivation of \isi{*\textit{h\textsubscript{2}e}-conjugation} presents is to some extent corroborated by the derivation of Hitt.\ \textit{ḫi}-verbs\iw{Hitt}{-\textit{ḫi}} \is{\textit{ḫi}-conjugation} from adverbs, as described by \citet{Melchert2009}: \iwi{Hitt}{\textit{āppai}-} ‘to go back, be finished’ besides \iwi{Hitt}{\textit{āppa}} ‘after, (temporally) behind’, \iwi{Hitt}{\textit{parā}-} ‘to appear, come/go forth’ (cf.\  HLuw.\  \iwi{HLuw}{\textit{ARHA para}-} ‘go missing, be missing, lack’) from \textit{parā}\iw{Hitt}{\textit{parā} (adv.)} ‘forth, out’, \iwi{Hitt}{\textit{šanna}-} ‘to conceal’ from *\textit{šanna}-\iw{Hitt}{*\textit{šanna}- (adv.)} ‘isolated, separated’ (< \iwi{PIE}{*\textit{sn̥(H)o(-)}}, cf.\  \iwi{Hitt}{\textit{šannapi}} ‘in an isolated place’). An example of the latter is given in (\ref{ex:10}).\footnote{Further instances in \citealt[111--113]{Puhvel2017}. The derivation proposed by \citet[336--337]{Melchert2009} is more straightforward than through a verbal stem with a nasal infix (\citealt[159--160]{Oettinger1979}; \citealt[719]{Kloekhorst2008}) or with a nasal suffix \citep[116]{Puhvel2017}.}  

\begin{exe}
    \ex KBo 5.3. I 28\\ \gll \textit{n-at-mu-kan} \textit{mān} {\textit{šannatti} \iw{Hitt}{\textit{šanna}-}} \textit{n-at-mu} \textit{ŪL} \textit{mematti}\\
    \textsc{conj-3sg.acc.n-1sg.acc/dat-ptcl} if conceal.\textsc{2sg.prs.act} \textsc{conj-3sg.acc.n-1sg.acc/dat} \textsc{neg} tell.\textsc{2sg.prs.act} \\
    \glt `if you keep it secret from me, and do not tell me' \label{ex:10}
\end{exe}


The secondary verbal inflection of local adverbs resides ultimately in the use of adverbs as predicate in nominal sentences. 

The process underlying the formation of the *-\textit{eh\textsubscript{2}}-\isi{factitive} is strikingly parallel to the derivation of the \isi{stative} in *-\textit{eh\textsubscript{1}}-\iw{PIE}{*-\textit{éh\textsubscript{1}-/-h\textsubscript{1}}-} from a predicative adverb based on the instrumental sg.\ of a root noun,\is{root!noun} e.g. \iwi{PIE}{*\textit{h\textsubscript{1}rudʰ-éh\textsubscript{1}}} ‘with redness’ → \isi{stative} \iwi{PIE}{*\textit{h\textsubscript{1}rudʰ-éh\textsubscript{1}}- ‘to be red’}, verbalized through addition of the present suffix \is{*-\textit{i̯e/o}-present}  *-\textit{i̯é/ó}-\iw{PIE}{*-\textit{i̯e/o}-}\footnote{See \citealt[145--149]{Jasanoff2004}; also Merritt, this volume.} (see the introduction to this volume). In some languages, the \isi{stative} itself coexisted with the productive\is{productivity} formations of \isi{fientive}s in \iwi{PIE}{*-\textit{eh\textsubscript{1}-s}-} and \isi{inchoative}s in \iwi{PIE}{*-\textit{eh\textsubscript{1}-s\kinvbreve e/o}-}, or of periphrastic\is{periphrasis} \isi{factitive}s, type Lat.\ \textit{ārē-faciō} \iw{Latin}{\textit{ārēfaciō}} ‘to make dry’, beside \textit{ārē-scō} \iw{Latin}{\textit{ārēscō}} ‘to grow dry’. In addition, the \isi{stative} present stem became the basis of nominal forms in \iwi{PIE}{*-\textit{eh\textsubscript{1}-to}-}, \iwi{PIE}{*-\textit{eh\textsubscript{1}-ti}-}, etc. The thematic renewal of the *-\textit{eh\textsubscript{2}}- \isi{factitive} type into \iwi{PIE}{*-\textit{eh\textsubscript{2}-i̯e/o}-} would be similar to the replacement of the *-\textit{eh\textsubscript{1}}-\isi{stative} type\iw{PIE}{*-\textit{éh\textsubscript{1}-/-h\textsubscript{1}}-} by \iwi{PIE}{*-\textit{eh\textsubscript{1}-i̯e/o}-}, e.g. in Lat.\   \iwi{Latin}{\textit{rubeō}}, -\textit{ēre}, etc. Parallel to the formations issued from the \isi{stative}, secondary adjectives\is{derivation!stem-based} in \iwi{PIE}{*-\textit{to}-} and \iwi{PIE}{*\textit{-(e/o)nt}-} based directly on the abstract probably also existed, which were later grammaticalized as participles\is{participle} of the newer *-\textit{ah\textsubscript{2}}-\isi{factitive}s,\is{\textit{nēwaḫḫi}-type} to wit Lat.\  \iwi{Latin}{\textit{(re)nouātus}}, Hitt.\ \iwi{Hitt}{\textit{nēwaḫḫant}-}.  

\section{Concluding remarks}

In the first two sections of this paper, two alleged instances of denominative presents\is{denominative!verb} without an overt suffix have been falsified. This was made possible by resorting to more economic analogical processes. As an exemplary point, the theory of \isi{conversion} of an adjective into a verb has been challenged in the case of \iwi{PIE}{*\textit{g\textsuperscript{u̯}ih\textsubscript{3}-u̯ó}-} ‘living’ and \iw{PIE}{*\textit{g\textsuperscript{u̯}ih\textsubscript{3}-u̯é/ó}-} *\textit{g\textsuperscript{u̯}ih\textsubscript{3}-u̯é-ti} ‘lives’. This case suggests that the analogical process took place already in PIE. In the third section, the intriguing case of \isi{conversion} apparently featured by the \isi{factitive} *-\textit{eh\textsubscript{2}}-present type has been treated in some detail and was resolved by resorting to a further analogical process in PIE. The crucial fact lies in the prevalent derivation of this \isi{factitive} type from abstracts and individualizing derivatives, as still reflected by the Anatolian data. Furthermore, the later productive\is{productivity} development of *-\textit{eh\textsubscript{2}}-abstracts from thematic \isi{adnominal} bases provides the missing link for deriving the *\textit{néu̯-eh\textsubscript{2}}-type\iw{PIE}{*\textit{néu̯-e(-)h\textsubscript{2}}- ‘to make new’} from thematic adjectives. The reconstructed process also explains without further ado the original \isi{*\textit{h\textsubscript{2}e}-conjugation} of the *-\textit{eh\textsubscript{2}}-\isi{factitive} type, which had been never accounted for. As an alternative approach, which would effectively be the admission of failure, one would be forced to assume that the verbal suffix \iwi{PIE}{*-\textit{eh\textsubscript{2}}-} of the present type and the nominal suffix \iwi{PIE}{*-\textit{eh\textsubscript{2}}-} had nothing in common and were simply homophonous. It does not seem commendable to multiply homophonous PIE morphemes only by projecting back the data on a purely formal basis without paying attention to the functions of the derivatives. This amounts to giving up on any perspective for a solution with the excuse of ignorance. With respect to the semantic side, the main issue concerned explaining the \isi{factitive} meaning of the denominative\is{denominative!verb} *-\textit{eh\textsubscript{2}}-present type, which follows from the syntactic reconstruction of predicatively used\is{predicative!usage} verbal abstracts with nominal complements.  

\section*{Abbreviations}
\begin{tabularx}{.48\textwidth}{lQ}
acc.\ & accusative\\
act.\ & active\\
adj.\ & adjective\\
aor.\ & aorist\\
Arm.\ & Armenian\\
AV & Saṃhitās of the Atharvaveda\\
Av.\ & Avestan\\
CLuw.\ & Cuneiform Luwian\\
CToch.\ & Common Tocharian\\
dat.\ & dative\\
fem.\ & feminine\\
Gaul.\ & Gaulish\\
gen.\ & genitive\\
Gk.\ & Greek\\
Go.\ & Gothic\\
HLuw.\ & Hieroglyphic Luwian\\
Hitt.\ & Hittite\\
IE & Indo-European\\
instr.\ & instrumental\\
Lat.\ & Latin\\
Latv.\ & Latvian\\
Lesb.\ & Lesbian\\
Lith.\ & Lithuanian\\
Lyc.\ & Lycian\\
Lyd.\ & Lydian\\
\end{tabularx}
\begin{tabularx}{.48\textwidth}{lQ}
MH & Middle Hittite\\
MP & Middle Persian\\
masc.\ & masculine\\
mid.\ & middle\\
Myc.\ & Mycenaean\\
nom.\ & nominative\\
nt.\ & neuter\\
OAv.\ & Old Avestan\\
OCS & Old Church Slavonic\\
OH & Old Hittite\\
OHG	& Old High German\\
OIr.\ & Old Irish\\
OP & Old Persian\\
OPr.\ & Old Prussian\\
Oss.\ & Ossetic\\
PIE & Proto-Indo-European\\
pl.\ & plural\\
pres.\ & present\\
ptcp.\ & participle\\
RV & Saṃhitā of the R̥gveda\\
sg.\ & singular\\
Toch.\ & Tocharian\\
Ved.\ & Vedic\\
VS\ & Vājasaneyi-Saṃhitā\\
YAv.\ & Young Avestan\\
\end{tabularx}

{\sloppy\printbibliography[heading=subbibliography,notkeyword=this]}
\end{document}
















